%************************************************
\chapter{The Algebra of General Numbers: A Formulary}\label{app:gn-algebra}
%************************************************
% Appendix F (added 6 November 2026), included by book.tex after %************************************************
\chapter{Identity Cards of Matter}\label{app:identity-cards}
%************************************************

This appendix lists the identity cards $(r, \alpha, \phi)$ at the angle $\pb \approx 79.86^\circ$, of the 306 bodies of Chapter~\ref{ch:standard-model}: the electron, the proton and the neutron; then, in the order of the periodic table, each of the 118 atoms, followed by its common ions and by its bare nucleus. The masses are those of \Cref{sec:sm-cards}: atoms from the standard atomic weights, ions as $A - z\,m_e$, and nuclei as $A - Z\,m_e + B_e/c^2$. $Z$ is the atomic number and $z$ the charge number. For a neutral body $r = m$, $\alpha = 0$ and $\phi = 79.86^\circ$. For a charged body $r$ is $|z|\,m_S$ to many digits, and the angles are written with the lean $\delta$ of Remark~\ref{rem:lean}, which is far too small to be written as a plain decimal. The workbook \emph{Identity-Cards-data.xlsx} holds the same numbers as live formulas.

\bigskip
\begingroup\scriptsize\setlength{\extrarowheight}{0pt}\setlength{\tabcolsep}{3.2pt}\renewcommand{\arraystretch}{1.02}
\noindent\makebox[\textwidth][c]{%
\begin{tabular}{@{}r l r r r r r l l@{}}
\toprule
& \tableheadline{body} & $Z$ & $z$ & $m$ (u) & $m$ (kg) & $r$ (kg) & $\alpha$ (degree) & $\phi$ (degree) \\
\midrule
1 & electron &  & $-1$ & $0.0005485799091$ & $9.10938{\times}10^{-31}$ & $1.85921{\times}10^{-9}$ & $259.86 + 2.76342{\times}10^{-20}$ & $-180 - 2.76342{\times}10^{-20}$ \\
2 & proton &  & $+1$ & $1.007276467$ & $1.67262{\times}10^{-27}$ & $1.85921{\times}10^{-9}$ & $79.86 - 5.07406{\times}10^{-17}$ & $5.07406{\times}10^{-17}$ \\
3 & neutron &  & $0$ & $1.008664916$ & $1.67493{\times}10^{-27}$ & $1.67493{\times}10^{-27}$ & $0$ & $79.86$ \\
4 & H & 1 & $0$ & $1.008$ & $1.67382{\times}10^{-27}$ & $1.67382{\times}10^{-27}$ & $0$ & $79.86$ \\
5 & H$^{+}$ & 1 & $+1$ & $1.00745142$ & $1.67291{\times}10^{-27}$ & $1.85921{\times}10^{-9}$ & $79.86 - 5.07494{\times}10^{-17}$ & $5.07494{\times}10^{-17}$ \\
6 & H$^{-}$ & 1 & $-1$ & $1.00854858$ & $1.67473{\times}10^{-27}$ & $1.85921{\times}10^{-9}$ & $259.86 + 5.08046{\times}10^{-17}$ & $-180 - 5.08046{\times}10^{-17}$ \\
7 & H nucleus & 1 & $+1$ & $1.007451436$ & $1.67291{\times}10^{-27}$ & $1.85921{\times}10^{-9}$ & $79.86 - 5.07494{\times}10^{-17}$ & $5.07494{\times}10^{-17}$ \\
8 & He & 2 & $0$ & $4.002602$ & $6.64648{\times}10^{-27}$ & $6.64648{\times}10^{-27}$ & $0$ & $79.86$ \\
9 & He nucleus & 2 & $+2$ & $4.001504921$ & $6.64466{\times}10^{-27}$ & $3.71842{\times}10^{-9}$ & $79.86 - 1.00786{\times}10^{-16}$ & $1.00786{\times}10^{-16}$ \\
10 & Li & 3 & $0$ & $6.94$ & $1.15241{\times}10^{-26}$ & $1.15241{\times}10^{-26}$ & $0$ & $79.86$ \\
11 & Li$^{+}$ & 3 & $+1$ & $6.93945142$ & $1.15232{\times}10^{-26}$ & $1.85921{\times}10^{-9}$ & $79.86 - 3.49568{\times}10^{-16}$ & $3.49568{\times}10^{-16}$ \\
12 & Li nucleus & 3 & $+3$ & $6.938354474$ & $1.15214{\times}10^{-26}$ & $5.57763{\times}10^{-9}$ & $79.86 - 1.16504{\times}10^{-16}$ & $1.16504{\times}10^{-16}$ \\
13 & Be & 4 & $0$ & $9.0121831$ & $1.49651{\times}10^{-26}$ & $1.49651{\times}10^{-26}$ & $0$ & $79.86$ \\
14 & Be$^{2+}$ & 4 & $+2$ & $9.01108594$ & $1.49633{\times}10^{-26}$ & $3.71842{\times}10^{-9}$ & $79.86 - 2.26962{\times}10^{-16}$ & $2.26962{\times}10^{-16}$ \\
15 & Be nucleus & 4 & $+4$ & $9.009989206$ & $1.49614{\times}10^{-26}$ & $7.43684{\times}10^{-9}$ & $79.86 - 1.13467{\times}10^{-16}$ & $1.13467{\times}10^{-16}$ \\
16 & B & 5 & $0$ & $10.81$ & $1.79504{\times}10^{-26}$ & $1.79504{\times}10^{-26}$ & $0$ & $79.86$ \\
17 & B nucleus & 5 & $+5$ & $10.80725783$ & $1.79459{\times}10^{-26}$ & $9.29605{\times}10^{-9}$ & $79.86 - 1.08881{\times}10^{-16}$ & $1.08881{\times}10^{-16}$ \\
18 & C & 6 & $0$ & $12.011$ & $1.99447{\times}10^{-26}$ & $1.99447{\times}10^{-26}$ & $0$ & $79.86$ \\
19 & C nucleus & 6 & $+6$ & $12.00770964$ & $1.99393{\times}10^{-26}$ & $1.11553{\times}10^{-8}$ & $79.86 - 1.00813{\times}10^{-16}$ & $1.00813{\times}10^{-16}$ \\
20 & N & 7 & $0$ & $14.007$ & $2.32592{\times}10^{-26}$ & $2.32592{\times}10^{-26}$ & $0$ & $79.86$ \\
21 & N$^{3-}$ & 7 & $-3$ & $14.00864574$ & $2.32619{\times}10^{-26}$ & $5.57763{\times}10^{-9}$ & $259.86 + 2.35224{\times}10^{-16}$ & $-180 - 2.35224{\times}10^{-16}$ \\
22 & N nucleus & 7 & $+7$ & $14.00316156$ & $2.32528{\times}10^{-26}$ & $1.30145{\times}10^{-8}$ & $79.86 - 1.00771{\times}10^{-16}$ & $1.00771{\times}10^{-16}$ \\
23 & O & 8 & $0$ & $15.999$ & $2.65670{\times}10^{-26}$ & $2.65670{\times}10^{-26}$ & $0$ & $79.86$ \\
24 & O$^{2-}$ & 8 & $-2$ & $16.00009716$ & $2.65688{\times}10^{-26}$ & $3.71842{\times}10^{-9}$ & $259.86 + 4.02995{\times}10^{-16}$ & $-180 - 4.02995{\times}10^{-16}$ \\
25 & O nucleus & 8 & $+8$ & $15.99461359$ & $2.65597{\times}10^{-26}$ & $1.48737{\times}10^{-8}$ & $79.86 - 1.00714{\times}10^{-16}$ & $1.00714{\times}10^{-16}$ \\
26 & F & 9 & $0$ & $18.99840316$ & $3.15476{\times}10^{-26}$ & $3.15476{\times}10^{-26}$ & $0$ & $79.86$ \\
27 & F$^{-}$ & 9 & $-1$ & $18.99895174$ & $3.15485{\times}10^{-26}$ & $1.85921{\times}10^{-9}$ & $259.86 + 9.57053{\times}10^{-16}$ & $-180 - 9.57053{\times}10^{-16}$ \\
28 & F nucleus & 9 & $+9$ & $18.9934689$ & $3.15394{\times}10^{-26}$ & $1.67329{\times}10^{-8}$ & $79.86 - 1.06309{\times}10^{-16}$ & $1.06309{\times}10^{-16}$ \\
29 & Ne & 10 & $0$ & $20.1797$ & $3.35092{\times}10^{-26}$ & $3.35092{\times}10^{-26}$ & $0$ & $79.86$ \\
30 & Ne nucleus & 10 & $+10$ & $20.17421801$ & $3.35001{\times}10^{-26}$ & $1.85921{\times}10^{-8}$ & $79.86 - 1.01626{\times}10^{-16}$ & $1.01626{\times}10^{-16}$ \\
31 & Na & 11 & $0$ & $22.98976928$ & $3.81754{\times}10^{-26}$ & $3.81754{\times}10^{-26}$ & $0$ & $79.86$ \\
32 & Na$^{+}$ & 11 & $+1$ & $22.9892207$ & $3.81745{\times}10^{-26}$ & $1.85921{\times}10^{-9}$ & $79.86 - 1.15806{\times}10^{-15}$ & $1.15806{\times}10^{-15}$ \\
33 & Na nucleus & 11 & $+11$ & $22.98373968$ & $3.81654{\times}10^{-26}$ & $2.04513{\times}10^{-8}$ & $79.86 - 1.05253{\times}10^{-16}$ & $1.05253{\times}10^{-16}$ \\
34 & Mg & 12 & $0$ & $24.305$ & $4.03594{\times}10^{-26}$ & $4.03594{\times}10^{-26}$ & $0$ & $79.86$ \\
35 & Mg$^{2+}$ & 12 & $+2$ & $24.30390284$ & $4.03576{\times}10^{-26}$ & $3.71842{\times}10^{-9}$ & $79.86 - 6.12142{\times}10^{-16}$ & $6.12142{\times}10^{-16}$ \\
36 & Mg nucleus & 12 & $+12$ & $24.29842292$ & $4.03485{\times}10^{-26}$ & $2.23105{\times}10^{-8}$ & $79.86 - 1.02001{\times}10^{-16}$ & $1.02001{\times}10^{-16}$ \\
\bottomrule
\end{tabular}}
\endgroup
\clearpage
\begingroup\scriptsize\setlength{\extrarowheight}{0pt}\setlength{\tabcolsep}{3.2pt}\renewcommand{\arraystretch}{1.02}
\noindent\makebox[\textwidth][c]{%
\begin{tabular}{@{}r l r r r r r l l@{}}
\toprule
& \tableheadline{body} & $Z$ & $z$ & $m$ (u) & $m$ (kg) & $r$ (kg) & $\alpha$ (degree) & $\phi$ (degree) \\
\midrule
37 & Al & 13 & $0$ & $26.9815384$ & $4.48039{\times}10^{-26}$ & $4.48039{\times}10^{-26}$ & $0$ & $79.86$ \\
38 & Al$^{3+}$ & 13 & $+3$ & $26.97989266$ & $4.48012{\times}10^{-26}$ & $5.57763{\times}10^{-9}$ & $79.86 - 4.53028{\times}10^{-16}$ & $4.53028{\times}10^{-16}$ \\
39 & Al nucleus & 13 & $+13$ & $26.97441399$ & $4.47921{\times}10^{-26}$ & $2.41697{\times}10^{-8}$ & $79.86 - 1.04524{\times}10^{-16}$ & $1.04524{\times}10^{-16}$ \\
40 & Si & 14 & $0$ & $28.085$ & $4.66362{\times}10^{-26}$ & $4.66362{\times}10^{-26}$ & $0$ & $79.86$ \\
41 & Si nucleus & 14 & $+14$ & $28.07732839$ & $4.66235{\times}10^{-26}$ & $2.60289{\times}10^{-8}$ & $79.86 - 1.01026{\times}10^{-16}$ & $1.01026{\times}10^{-16}$ \\
42 & P & 15 & $0$ & $30.973762$ & $5.14331{\times}10^{-26}$ & $5.14331{\times}10^{-26}$ & $0$ & $79.86$ \\
43 & P$^{3-}$ & 15 & $-3$ & $30.97540774$ & $5.14359{\times}10^{-26}$ & $5.57763{\times}10^{-9}$ & $259.86 + 5.20118{\times}10^{-16}$ & $-180 - 5.20118{\times}10^{-16}$ \\
44 & P nucleus & 15 & $+15$ & $30.96554333$ & $5.14195{\times}10^{-26}$ & $2.78881{\times}10^{-8}$ & $79.86 - 1.03991{\times}10^{-16}$ & $1.03991{\times}10^{-16}$ \\
45 & S & 16 & $0$ & $32.06$ & $5.32369{\times}10^{-26}$ & $5.32369{\times}10^{-26}$ & $0$ & $79.86$ \\
46 & S$^{2-}$ & 16 & $-2$ & $32.06109716$ & $5.32387{\times}10^{-26}$ & $3.71842{\times}10^{-9}$ & $259.86 + 8.07523{\times}10^{-16}$ & $-180 - 8.07523{\times}10^{-16}$ \\
47 & S nucleus & 16 & $+16$ & $32.05123443$ & $5.32223{\times}10^{-26}$ & $2.97473{\times}10^{-8}$ & $79.86 - 1.00909{\times}10^{-16}$ & $1.00909{\times}10^{-16}$ \\
48 & Cl & 17 & $0$ & $35.45$ & $5.88661{\times}10^{-26}$ & $5.88661{\times}10^{-26}$ & $0$ & $79.86$ \\
49 & Cl$^{-}$ & 17 & $-1$ & $35.45054858$ & $5.88670{\times}10^{-26}$ & $1.85921{\times}10^{-9}$ & $259.86 + 1.78579{\times}10^{-15}$ & $-180 - 1.78579{\times}10^{-15}$ \\
50 & Cl nucleus & 17 & $+17$ & $35.44068767$ & $5.88506{\times}10^{-26}$ & $3.16066{\times}10^{-8}$ & $79.86 - 1.05017{\times}10^{-16}$ & $1.05017{\times}10^{-16}$ \\
51 & Ar & 18 & $0$ & $39.95$ & $6.63385{\times}10^{-26}$ & $6.63385{\times}10^{-26}$ & $0$ & $79.86$ \\
52 & Ar nucleus & 18 & $+18$ & $39.94014107$ & $6.63222{\times}10^{-26}$ & $3.34658{\times}10^{-8}$ & $79.86 - 1.11775{\times}10^{-16}$ & $1.11775{\times}10^{-16}$ \\
53 & K & 19 & $0$ & $39.0983$ & $6.49243{\times}10^{-26}$ & $6.49243{\times}10^{-26}$ & $0$ & $79.86$ \\
54 & K$^{+}$ & 19 & $+1$ & $39.09775142$ & $6.49233{\times}10^{-26}$ & $1.85921{\times}10^{-9}$ & $79.86 - 1.96951{\times}10^{-15}$ & $1.96951{\times}10^{-15}$ \\
55 & K nucleus & 19 & $+19$ & $39.08789464$ & $6.49070{\times}10^{-26}$ & $3.53250{\times}10^{-8}$ & $79.86 - 1.03632{\times}10^{-16}$ & $1.03632{\times}10^{-16}$ \\
56 & Ca & 20 & $0$ & $40.078$ & $6.65511{\times}10^{-26}$ & $6.65511{\times}10^{-26}$ & $0$ & $79.86$ \\
57 & Ca$^{2+}$ & 20 & $+2$ & $40.07690284$ & $6.65493{\times}10^{-26}$ & $3.71842{\times}10^{-9}$ & $79.86 - 1.00942{\times}10^{-15}$ & $1.00942{\times}10^{-15}$ \\
58 & Ca nucleus & 20 & $+20$ & $40.06704836$ & $6.65329{\times}10^{-26}$ & $3.71842{\times}10^{-8}$ & $79.86 - 1.00917{\times}10^{-16}$ & $1.00917{\times}10^{-16}$ \\
59 & Sc & 21 & $0$ & $44.955907$ & $7.46510{\times}10^{-26}$ & $7.46510{\times}10^{-26}$ & $0$ & $79.86$ \\
60 & Sc$^{3+}$ & 21 & $+3$ & $44.95426126$ & $7.46483{\times}10^{-26}$ & $5.57763{\times}10^{-9}$ & $79.86 - 7.54842{\times}10^{-16}$ & $7.54842{\times}10^{-16}$ \\
61 & Sc nucleus & 21 & $+21$ & $44.94440925$ & $7.46319{\times}10^{-26}$ & $3.90434{\times}10^{-8}$ & $79.86 - 1.07811{\times}10^{-16}$ & $1.07811{\times}10^{-16}$ \\
62 & Ti & 22 & $0$ & $47.867$ & $7.94850{\times}10^{-26}$ & $7.94850{\times}10^{-26}$ & $0$ & $79.86$ \\
63 & Ti nucleus & 22 & $+22$ & $47.85495631$ & $7.94650{\times}10^{-26}$ & $4.09026{\times}10^{-8}$ & $79.86 - 1.09575{\times}10^{-16}$ & $1.09575{\times}10^{-16}$ \\
64 & V & 23 & $0$ & $50.9415$ & $8.45904{\times}10^{-26}$ & $8.45904{\times}10^{-26}$ & $0$ & $79.86$ \\
65 & V nucleus & 23 & $+23$ & $50.92891055$ & $8.45694{\times}10^{-26}$ & $4.27618{\times}10^{-8}$ & $79.86 - 1.11543{\times}10^{-16}$ & $1.11543{\times}10^{-16}$ \\
66 & Cr & 24 & $0$ & $51.9961$ & $8.63416{\times}10^{-26}$ & $8.63416{\times}10^{-26}$ & $0$ & $79.86$ \\
67 & Cr$^{2+}$ & 24 & $+2$ & $51.99500284$ & $8.63397{\times}10^{-26}$ & $3.71842{\times}10^{-9}$ & $79.86 - 1.30960{\times}10^{-15}$ & $1.30960{\times}10^{-15}$ \\
68 & Cr$^{3+}$ & 24 & $+3$ & $51.99445426$ & $8.63388{\times}10^{-26}$ & $5.57763{\times}10^{-9}$ & $79.86 - 8.73056{\times}10^{-16}$ & $8.73056{\times}10^{-16}$ \\
69 & Cr nucleus & 24 & $+24$ & $51.98296496$ & $8.63197{\times}10^{-26}$ & $4.46210{\times}10^{-8}$ & $79.86 - 1.09108{\times}10^{-16}$ & $1.09108{\times}10^{-16}$ \\
70 & Mn & 25 & $0$ & $54.938043$ & $9.12268{\times}10^{-26}$ & $9.12268{\times}10^{-26}$ & $0$ & $79.86$ \\
71 & Mn$^{2+}$ & 25 & $+2$ & $54.93694584$ & $9.12249{\times}10^{-26}$ & $3.71842{\times}10^{-9}$ & $79.86 - 1.38370{\times}10^{-15}$ & $1.38370{\times}10^{-15}$ \\
72 & Mn$^{3+}$ & 25 & $+3$ & $54.93639726$ & $9.12240{\times}10^{-26}$ & $5.57763{\times}10^{-9}$ & $79.86 - 9.22455{\times}10^{-16}$ & $9.22455{\times}10^{-16}$ \\
73 & Mn$^{4+}$ & 25 & $+4$ & $54.93584868$ & $9.12231{\times}10^{-26}$ & $7.43684{\times}10^{-9}$ & $79.86 - 6.91835{\times}10^{-16}$ & $6.91835{\times}10^{-16}$ \\
74 & Mn nucleus & 25 & $+25$ & $54.92436255$ & $9.12040{\times}10^{-26}$ & $4.64802{\times}10^{-8}$ & $79.86 - 1.10670{\times}10^{-16}$ & $1.10670{\times}10^{-16}$ \\
75 & Fe & 26 & $0$ & $55.845$ & $9.27328{\times}10^{-26}$ & $9.27328{\times}10^{-26}$ & $0$ & $79.86$ \\
76 & Fe$^{2+}$ & 26 & $+2$ & $55.84390284$ & $9.27310{\times}10^{-26}$ & $3.71842{\times}10^{-9}$ & $79.86 - 1.40654{\times}10^{-15}$ & $1.40654{\times}10^{-15}$ \\
77 & Fe$^{3+}$ & 26 & $+3$ & $55.84335426$ & $9.27301{\times}10^{-26}$ & $5.57763{\times}10^{-9}$ & $79.86 - 9.37685{\times}10^{-16}$ & $9.37685{\times}10^{-16}$ \\
78 & Fe nucleus & 26 & $+26$ & $55.83077432$ & $9.27092{\times}10^{-26}$ & $4.83394{\times}10^{-8}$ & $79.86 - 1.08170{\times}10^{-16}$ & $1.08170{\times}10^{-16}$ \\
79 & Co & 27 & $0$ & $58.933194$ & $9.78609{\times}10^{-26}$ & $9.78609{\times}10^{-26}$ & $0$ & $79.86$ \\
80 & Co$^{2+}$ & 27 & $+2$ & $58.93209684$ & $9.78590{\times}10^{-26}$ & $3.71842{\times}10^{-9}$ & $79.86 - 1.48432{\times}10^{-15}$ & $1.48432{\times}10^{-15}$ \\
81 & Co$^{3+}$ & 27 & $+3$ & $58.93154826$ & $9.78581{\times}10^{-26}$ & $5.57763{\times}10^{-9}$ & $79.86 - 9.89539{\times}10^{-16}$ & $9.89539{\times}10^{-16}$ \\
82 & Co nucleus & 27 & $+27$ & $58.91842328$ & $9.78363{\times}10^{-26}$ & $5.01986{\times}10^{-8}$ & $79.86 - 1.09924{\times}10^{-16}$ & $1.09924{\times}10^{-16}$ \\
83 & Ni & 28 & $0$ & $58.6934$ & $9.74627{\times}10^{-26}$ & $9.74627{\times}10^{-26}$ & $0$ & $79.86$ \\
84 & Ni nucleus & 28 & $+28$ & $58.67808442$ & $9.74373{\times}10^{-26}$ & $5.20579{\times}10^{-8}$ & $79.86 - 1.05566{\times}10^{-16}$ & $1.05566{\times}10^{-16}$ \\
85 & Cu & 29 & $0$ & $63.546$ & $1.05521{\times}10^{-25}$ & $1.05521{\times}10^{-25}$ & $0$ & $79.86$ \\
86 & Cu$^{+}$ & 29 & $+1$ & $63.54545142$ & $1.05520{\times}10^{-25}$ & $1.85921{\times}10^{-9}$ & $79.86 - 3.20104{\times}10^{-15}$ & $3.20104{\times}10^{-15}$ \\
87 & Cu$^{2+}$ & 29 & $+2$ & $63.54490284$ & $1.05519{\times}10^{-25}$ & $3.71842{\times}10^{-9}$ & $79.86 - 1.60051{\times}10^{-15}$ & $1.60051{\times}10^{-15}$ \\
88 & Cu nucleus & 29 & $+29$ & $63.53013976$ & $1.05494{\times}10^{-25}$ & $5.39171{\times}10^{-8}$ & $79.86 - 1.10354{\times}10^{-16}$ & $1.10354{\times}10^{-16}$ \\
89 & Zn & 30 & $0$ & $65.38$ & $1.08566{\times}10^{-25}$ & $1.08566{\times}10^{-25}$ & $0$ & $79.86$ \\
90 & Zn$^{2+}$ & 30 & $+2$ & $65.37890284$ & $1.08564{\times}10^{-25}$ & $3.71842{\times}10^{-9}$ & $79.86 - 1.64670{\times}10^{-15}$ & $1.64670{\times}10^{-15}$ \\
91 & Zn nucleus & 30 & $+30$ & $65.3635953$ & $1.08539{\times}10^{-25}$ & $5.57763{\times}10^{-8}$ & $79.86 - 1.09754{\times}10^{-16}$ & $1.09754{\times}10^{-16}$ \\
92 & Ga & 31 & $0$ & $69.723$ & $1.15778{\times}10^{-25}$ & $1.15778{\times}10^{-25}$ & $0$ & $79.86$ \\
93 & Ga$^{3+}$ & 31 & $+3$ & $69.72135426$ & $1.15775{\times}10^{-25}$ & $5.57763{\times}10^{-9}$ & $79.86 - 1.17071{\times}10^{-15}$ & $1.17071{\times}10^{-15}$ \\
94 & Ga nucleus & 31 & $+31$ & $69.70605103$ & $1.15750{\times}10^{-25}$ & $5.76355{\times}10^{-8}$ & $79.86 - 1.13270{\times}10^{-16}$ & $1.13270{\times}10^{-16}$ \\
95 & Ge & 32 & $0$ & $72.63$ & $1.20605{\times}10^{-25}$ & $1.20605{\times}10^{-25}$ & $0$ & $79.86$ \\
96 & Ge nucleus & 32 & $+32$ & $72.61250696$ & $1.20576{\times}10^{-25}$ & $5.94947{\times}10^{-8}$ & $79.86 - 1.14306{\times}10^{-16}$ & $1.14306{\times}10^{-16}$ \\
\bottomrule
\end{tabular}}
\endgroup
\clearpage
\begingroup\scriptsize\setlength{\extrarowheight}{0pt}\setlength{\tabcolsep}{3.2pt}\renewcommand{\arraystretch}{1.02}
\noindent\makebox[\textwidth][c]{%
\begin{tabular}{@{}r l r r r r r l l@{}}
\toprule
& \tableheadline{body} & $Z$ & $z$ & $m$ (u) & $m$ (kg) & $r$ (kg) & $\alpha$ (degree) & $\phi$ (degree) \\
\midrule
97 & As & 33 & $0$ & $74.921595$ & $1.24410{\times}10^{-25}$ & $1.24410{\times}10^{-25}$ & $0$ & $79.86$ \\
98 & As$^{3-}$ & 33 & $-3$ & $74.92324074$ & $1.24413{\times}10^{-25}$ & $5.57763{\times}10^{-9}$ & $259.86 + 1.25806{\times}10^{-15}$ & $-180 - 1.25806{\times}10^{-15}$ \\
99 & As nucleus & 33 & $+33$ & $74.90355809$ & $1.24380{\times}10^{-25}$ & $6.13539{\times}10^{-8}$ & $79.86 - 1.14339{\times}10^{-16}$ & $1.14339{\times}10^{-16}$ \\
100 & Se & 34 & $0$ & $78.971$ & $1.31134{\times}10^{-25}$ & $1.31134{\times}10^{-25}$ & $0$ & $79.86$ \\
101 & Se$^{2-}$ & 34 & $-2$ & $78.97209716$ & $1.31136{\times}10^{-25}$ & $3.71842{\times}10^{-9}$ & $259.86 + 1.98907{\times}10^{-15}$ & $-180 - 1.98907{\times}10^{-15}$ \\
102 & Se nucleus & 34 & $+34$ & $78.95241943$ & $1.31104{\times}10^{-25}$ & $6.32131{\times}10^{-8}$ & $79.86 - 1.16975{\times}10^{-16}$ & $1.16975{\times}10^{-16}$ \\
103 & Br & 35 & $0$ & $79.904$ & $1.32684{\times}10^{-25}$ & $1.32684{\times}10^{-25}$ & $0$ & $79.86$ \\
104 & Br$^{-}$ & 35 & $-1$ & $79.90454858$ & $1.32685{\times}10^{-25}$ & $1.85921{\times}10^{-9}$ & $259.86 + 4.02511{\times}10^{-15}$ & $-180 - 4.02511{\times}10^{-15}$ \\
105 & Br nucleus & 35 & $+35$ & $79.88487598$ & $1.32652{\times}10^{-25}$ & $6.50723{\times}10^{-8}$ & $79.86 - 1.14975{\times}10^{-16}$ & $1.14975{\times}10^{-16}$ \\
106 & Kr & 36 & $0$ & $83.798$ & $1.39150{\times}10^{-25}$ & $1.39150{\times}10^{-25}$ & $0$ & $79.86$ \\
107 & Kr nucleus & 36 & $+36$ & $83.77833274$ & $1.39117{\times}10^{-25}$ & $6.69315{\times}10^{-8}$ & $79.86 - 1.17229{\times}10^{-16}$ & $1.17229{\times}10^{-16}$ \\
108 & Rb & 37 & $0$ & $85.4678$ & $1.41923{\times}10^{-25}$ & $1.41923{\times}10^{-25}$ & $0$ & $79.86$ \\
109 & Rb$^{+}$ & 37 & $+1$ & $85.46725142$ & $1.41922{\times}10^{-25}$ & $1.85921{\times}10^{-9}$ & $79.86 - 4.30533{\times}10^{-15}$ & $4.30533{\times}10^{-15}$ \\
110 & Rb nucleus & 37 & $+37$ & $85.44758972$ & $1.41889{\times}10^{-25}$ & $6.87907{\times}10^{-8}$ & $79.86 - 1.16333{\times}10^{-16}$ & $1.16333{\times}10^{-16}$ \\
111 & Sr & 38 & $0$ & $87.62$ & $1.45496{\times}10^{-25}$ & $1.45496{\times}10^{-25}$ & $0$ & $79.86$ \\
112 & Sr$^{2+}$ & 38 & $+2$ & $87.61890284$ & $1.45495{\times}10^{-25}$ & $3.71842{\times}10^{-9}$ & $79.86 - 2.20686{\times}10^{-15}$ & $2.20686{\times}10^{-15}$ \\
113 & Sr nucleus & 38 & $+38$ & $87.59924691$ & $1.45462{\times}10^{-25}$ & $7.06499{\times}10^{-8}$ & $79.86 - 1.16124{\times}10^{-16}$ & $1.16124{\times}10^{-16}$ \\
114 & Y & 39 & $0$ & $88.905838$ & $1.47632{\times}10^{-25}$ & $1.47632{\times}10^{-25}$ & $0$ & $79.86$ \\
115 & Y$^{3+}$ & 39 & $+3$ & $88.90419226$ & $1.47629{\times}10^{-25}$ & $5.57763{\times}10^{-9}$ & $79.86 - 1.49282{\times}10^{-15}$ & $1.49282{\times}10^{-15}$ \\
116 & Y nucleus & 39 & $+39$ & $88.88454232$ & $1.47596{\times}10^{-25}$ & $7.25092{\times}10^{-8}$ & $79.86 - 1.14807{\times}10^{-16}$ & $1.14807{\times}10^{-16}$ \\
117 & Zr & 40 & $0$ & $91.222$ & $1.51478{\times}10^{-25}$ & $1.51478{\times}10^{-25}$ & $0$ & $79.86$ \\
118 & Zr nucleus & 40 & $+40$ & $91.20016196$ & $1.51441{\times}10^{-25}$ & $7.43684{\times}10^{-8}$ & $79.86 - 1.14853{\times}10^{-16}$ & $1.14853{\times}10^{-16}$ \\
119 & Nb & 41 & $0$ & $92.90637$ & $1.54275{\times}10^{-25}$ & $1.54275{\times}10^{-25}$ & $0$ & $79.86$ \\
120 & Nb nucleus & 41 & $+41$ & $92.88398982$ & $1.54237{\times}10^{-25}$ & $7.62276{\times}10^{-8}$ & $79.86 - 1.14120{\times}10^{-16}$ & $1.14120{\times}10^{-16}$ \\
121 & Mo & 42 & $0$ & $95.95$ & $1.59329{\times}10^{-25}$ & $1.59329{\times}10^{-25}$ & $0$ & $79.86$ \\
122 & Mo nucleus & 42 & $+42$ & $95.92707791$ & $1.59291{\times}10^{-25}$ & $7.80868{\times}10^{-8}$ & $79.86 - 1.15053{\times}10^{-16}$ & $1.15053{\times}10^{-16}$ \\
123 & Tc & 43 & $0$ & $97.0$ & $1.61072{\times}10^{-25}$ & $1.61072{\times}10^{-25}$ & $0$ & $79.86$ \\
124 & Tc nucleus & 43 & $+43$ & $96.97653624$ & $1.61033{\times}10^{-25}$ & $7.99460{\times}10^{-8}$ & $79.86 - 1.13607{\times}10^{-16}$ & $1.13607{\times}10^{-16}$ \\
125 & Ru & 44 & $0$ & $101.07$ & $1.67831{\times}10^{-25}$ & $1.67831{\times}10^{-25}$ & $0$ & $79.86$ \\
126 & Ru nucleus & 44 & $+44$ & $101.0459948$ & $1.67791{\times}10^{-25}$ & $8.18052{\times}10^{-8}$ & $79.86 - 1.15684{\times}10^{-16}$ & $1.15684{\times}10^{-16}$ \\
127 & Rh & 45 & $0$ & $102.90549$ & $1.70879{\times}10^{-25}$ & $1.70879{\times}10^{-25}$ & $0$ & $79.86$ \\
128 & Rh nucleus & 45 & $+45$ & $102.8809436$ & $1.70838{\times}10^{-25}$ & $8.36644{\times}10^{-8}$ & $79.86 - 1.15167{\times}10^{-16}$ & $1.15167{\times}10^{-16}$ \\
129 & Pd & 46 & $0$ & $106.42$ & $1.76715{\times}10^{-25}$ & $1.76715{\times}10^{-25}$ & $0$ & $79.86$ \\
130 & Pd nucleus & 46 & $+46$ & $106.3949126$ & $1.76673{\times}10^{-25}$ & $8.55236{\times}10^{-8}$ & $79.86 - 1.16512{\times}10^{-16}$ & $1.16512{\times}10^{-16}$ \\
131 & Ag & 47 & $0$ & $107.8682$ & $1.79119{\times}10^{-25}$ & $1.79119{\times}10^{-25}$ & $0$ & $79.86$ \\
132 & Ag$^{+}$ & 47 & $+1$ & $107.8676514$ & $1.79118{\times}10^{-25}$ & $1.85921{\times}10^{-9}$ & $79.86 - 5.43373{\times}10^{-15}$ & $5.43373{\times}10^{-15}$ \\
133 & Ag nucleus & 47 & $+47$ & $107.8425719$ & $1.79077{\times}10^{-25}$ & $8.73828{\times}10^{-8}$ & $79.86 - 1.15584{\times}10^{-16}$ & $1.15584{\times}10^{-16}$ \\
134 & Cd & 48 & $0$ & $112.414$ & $1.86668{\times}10^{-25}$ & $1.86668{\times}10^{-25}$ & $0$ & $79.86$ \\
135 & Cd$^{2+}$ & 48 & $+2$ & $112.4129028$ & $1.86666{\times}10^{-25}$ & $3.71842{\times}10^{-9}$ & $79.86 - 2.83134{\times}10^{-15}$ & $2.83134{\times}10^{-15}$ \\
136 & Cd nucleus & 48 & $+48$ & $112.3878314$ & $1.86624{\times}10^{-25}$ & $8.92420{\times}10^{-8}$ & $79.86 - 1.17946{\times}10^{-16}$ & $1.17946{\times}10^{-16}$ \\
137 & In & 49 & $0$ & $114.818$ & $1.90660{\times}10^{-25}$ & $1.90660{\times}10^{-25}$ & $0$ & $79.86$ \\
138 & In$^{3+}$ & 49 & $+3$ & $114.8163543$ & $1.90657{\times}10^{-25}$ & $5.57763{\times}10^{-9}$ & $79.86 - 1.92792{\times}10^{-15}$ & $1.92792{\times}10^{-15}$ \\
139 & In nucleus & 49 & $+49$ & $114.7912912$ & $1.90615{\times}10^{-25}$ & $9.11012{\times}10^{-8}$ & $79.86 - 1.18010{\times}10^{-16}$ & $1.18010{\times}10^{-16}$ \\
140 & Sn & 50 & $0$ & $118.71$ & $1.97123{\times}10^{-25}$ & $1.97123{\times}10^{-25}$ & $0$ & $79.86$ \\
141 & Sn$^{2+}$ & 50 & $+2$ & $118.7089028$ & $1.97121{\times}10^{-25}$ & $3.71842{\times}10^{-9}$ & $79.86 - 2.98992{\times}10^{-15}$ & $2.98992{\times}10^{-15}$ \\
142 & Sn$^{4+}$ & 50 & $+4$ & $118.7078057$ & $1.97119{\times}10^{-25}$ & $7.43684{\times}10^{-9}$ & $79.86 - 1.49495{\times}10^{-15}$ & $1.49495{\times}10^{-15}$ \\
143 & Sn nucleus & 50 & $+50$ & $118.6827512$ & $1.97077{\times}10^{-25}$ & $9.29605{\times}10^{-8}$ & $79.86 - 1.19571{\times}10^{-16}$ & $1.19571{\times}10^{-16}$ \\
144 & Sb & 51 & $0$ & $121.76$ & $2.02187{\times}10^{-25}$ & $2.02187{\times}10^{-25}$ & $0$ & $79.86$ \\
145 & Sb nucleus & 51 & $+51$ & $121.7322115$ & $2.02141{\times}10^{-25}$ & $9.48197{\times}10^{-8}$ & $79.86 - 1.20238{\times}10^{-16}$ & $1.20238{\times}10^{-16}$ \\
146 & Te & 52 & $0$ & $127.6$ & $2.11885{\times}10^{-25}$ & $2.11885{\times}10^{-25}$ & $0$ & $79.86$ \\
147 & Te$^{2-}$ & 52 & $-2$ & $127.6010972$ & $2.11887{\times}10^{-25}$ & $3.71842{\times}10^{-9}$ & $259.86 + 3.21389{\times}10^{-15}$ & $-180 - 3.21389{\times}10^{-15}$ \\
148 & Te nucleus & 52 & $+52$ & $127.5716721$ & $2.11838{\times}10^{-25}$ & $9.66789{\times}10^{-8}$ & $79.86 - 1.23583{\times}10^{-16}$ & $1.23583{\times}10^{-16}$ \\
149 & I & 53 & $0$ & $126.90447$ & $2.10730{\times}10^{-25}$ & $2.10730{\times}10^{-25}$ & $0$ & $79.86$ \\
150 & I$^{-}$ & 53 & $-1$ & $126.9050186$ & $2.10731{\times}10^{-25}$ & $1.85921{\times}10^{-9}$ & $259.86 + 6.39271{\times}10^{-15}$ & $-180 - 6.39271{\times}10^{-15}$ \\
151 & I nucleus & 53 & $+53$ & $126.8756029$ & $2.10682{\times}10^{-25}$ & $9.85381{\times}10^{-8}$ & $79.86 - 1.20589{\times}10^{-16}$ & $1.20589{\times}10^{-16}$ \\
152 & Xe & 54 & $0$ & $131.293$ & $2.18017{\times}10^{-25}$ & $2.18017{\times}10^{-25}$ & $0$ & $79.86$ \\
153 & Xe nucleus & 54 & $+54$ & $131.2635939$ & $2.17968{\times}10^{-25}$ & $1.00397{\times}10^{-7}$ & $79.86 - 1.22450{\times}10^{-16}$ & $1.22450{\times}10^{-16}$ \\
154 & Cs & 55 & $0$ & $132.905452$ & $2.20695{\times}10^{-25}$ & $2.20695{\times}10^{-25}$ & $0$ & $79.86$ \\
155 & Cs$^{+}$ & 55 & $+1$ & $132.9049034$ & $2.20694{\times}10^{-25}$ & $1.85921{\times}10^{-9}$ & $79.86 - 6.69495{\times}10^{-15}$ & $6.69495{\times}10^{-15}$ \\
156 & Cs nucleus & 55 & $+55$ & $132.8755072$ & $2.20645{\times}10^{-25}$ & $1.02257{\times}10^{-7}$ & $79.86 - 1.21699{\times}10^{-16}$ & $1.21699{\times}10^{-16}$ \\
\bottomrule
\end{tabular}}
\endgroup
\clearpage
\begingroup\scriptsize\setlength{\extrarowheight}{0pt}\setlength{\tabcolsep}{3.2pt}\renewcommand{\arraystretch}{1.02}
\noindent\makebox[\textwidth][c]{%
\begin{tabular}{@{}r l r r r r r l l@{}}
\toprule
& \tableheadline{body} & $Z$ & $z$ & $m$ (u) & $m$ (kg) & $r$ (kg) & $\alpha$ (degree) & $\phi$ (degree) \\
\midrule
157 & Ba & 56 & $0$ & $137.327$ & $2.28037{\times}10^{-25}$ & $2.28037{\times}10^{-25}$ & $0$ & $79.86$ \\
158 & Ba$^{2+}$ & 56 & $+2$ & $137.3259028$ & $2.28035{\times}10^{-25}$ & $3.71842{\times}10^{-9}$ & $79.86 - 3.45883{\times}10^{-15}$ & $3.45883{\times}10^{-15}$ \\
159 & Ba nucleus & 56 & $+56$ & $137.2965169$ & $2.27986{\times}10^{-25}$ & $1.04116{\times}10^{-7}$ & $79.86 - 1.23503{\times}10^{-16}$ & $1.23503{\times}10^{-16}$ \\
160 & La & 57 & $0$ & $138.90547$ & $2.30658{\times}10^{-25}$ & $2.30658{\times}10^{-25}$ & $0$ & $79.86$ \\
161 & La$^{3+}$ & 57 & $+3$ & $138.9038243$ & $2.30655{\times}10^{-25}$ & $5.57763{\times}10^{-9}$ & $79.86 - 2.33238{\times}10^{-15}$ & $2.33238{\times}10^{-15}$ \\
162 & La nucleus & 57 & $+57$ & $138.8744488$ & $2.30606{\times}10^{-25}$ & $1.05975{\times}10^{-7}$ & $79.86 - 1.22731{\times}10^{-16}$ & $1.22731{\times}10^{-16}$ \\
163 & Ce & 58 & $0$ & $140.116$ & $2.32668{\times}10^{-25}$ & $2.32668{\times}10^{-25}$ & $0$ & $79.86$ \\
164 & Ce$^{3+}$ & 58 & $+3$ & $140.1143543$ & $2.32665{\times}10^{-25}$ & $5.57763{\times}10^{-9}$ & $79.86 - 2.35271{\times}10^{-15}$ & $2.35271{\times}10^{-15}$ \\
165 & Ce nucleus & 58 & $+58$ & $140.084441$ & $2.32616{\times}10^{-25}$ & $1.07834{\times}10^{-7}$ & $79.86 - 1.21666{\times}10^{-16}$ & $1.21666{\times}10^{-16}$ \\
166 & Pr & 59 & $0$ & $140.90766$ & $2.33983{\times}10^{-25}$ & $2.33983{\times}10^{-25}$ & $0$ & $79.86$ \\
167 & Pr$^{3+}$ & 59 & $+3$ & $140.9060143$ & $2.33980{\times}10^{-25}$ & $5.57763{\times}10^{-9}$ & $79.86 - 2.36600{\times}10^{-15}$ & $2.36600{\times}10^{-15}$ \\
168 & Pr nucleus & 59 & $+59$ & $140.8755634$ & $2.33929{\times}10^{-25}$ & $1.09693{\times}10^{-7}$ & $79.86 - 1.20279{\times}10^{-16}$ & $1.20279{\times}10^{-16}$ \\
169 & Nd & 60 & $0$ & $144.242$ & $2.39519{\times}10^{-25}$ & $2.39519{\times}10^{-25}$ & $0$ & $79.86$ \\
170 & Nd$^{3+}$ & 60 & $+3$ & $144.2403543$ & $2.39517{\times}10^{-25}$ & $5.57763{\times}10^{-9}$ & $79.86 - 2.42199{\times}10^{-15}$ & $2.42199{\times}10^{-15}$ \\
171 & Nd nucleus & 60 & $+60$ & $144.2093661$ & $2.39465{\times}10^{-25}$ & $1.11553{\times}10^{-7}$ & $79.86 - 1.21073{\times}10^{-16}$ & $1.21073{\times}10^{-16}$ \\
172 & Pm & 61 & $0$ & $145.0$ & $2.40778{\times}10^{-25}$ & $2.40778{\times}10^{-25}$ & $0$ & $79.86$ \\
173 & Pm$^{3+}$ & 61 & $+3$ & $144.9983543$ & $2.40775{\times}10^{-25}$ & $5.57763{\times}10^{-9}$ & $79.86 - 2.43472{\times}10^{-15}$ & $2.43472{\times}10^{-15}$ \\
174 & Pm nucleus & 61 & $+61$ & $144.9668292$ & $2.40723{\times}10^{-25}$ & $1.13412{\times}10^{-7}$ & $79.86 - 1.19714{\times}10^{-16}$ & $1.19714{\times}10^{-16}$ \\
175 & Sm & 62 & $0$ & $150.36$ & $2.49679{\times}10^{-25}$ & $2.49679{\times}10^{-25}$ & $0$ & $79.86$ \\
176 & Sm$^{3+}$ & 62 & $+3$ & $150.3583543$ & $2.49676{\times}10^{-25}$ & $5.57763{\times}10^{-9}$ & $79.86 - 2.52472{\times}10^{-15}$ & $2.52472{\times}10^{-15}$ \\
177 & Sm nucleus & 62 & $+62$ & $150.3262925$ & $2.49623{\times}10^{-25}$ & $1.15271{\times}10^{-7}$ & $79.86 - 1.22138{\times}10^{-16}$ & $1.22138{\times}10^{-16}$ \\
178 & Eu & 63 & $0$ & $151.964$ & $2.52342{\times}10^{-25}$ & $2.52342{\times}10^{-25}$ & $0$ & $79.86$ \\
179 & Eu$^{3+}$ & 63 & $+3$ & $151.9623543$ & $2.52339{\times}10^{-25}$ & $5.57763{\times}10^{-9}$ & $79.86 - 2.55165{\times}10^{-15}$ & $2.55165{\times}10^{-15}$ \\
180 & Eu nucleus & 63 & $+63$ & $151.9297561$ & $2.52285{\times}10^{-25}$ & $1.17130{\times}10^{-7}$ & $79.86 - 1.21481{\times}10^{-16}$ & $1.21481{\times}10^{-16}$ \\
181 & Gd & 64 & $0$ & $157.249$ & $2.61118{\times}10^{-25}$ & $2.61118{\times}10^{-25}$ & $0$ & $79.86$ \\
182 & Gd$^{3+}$ & 64 & $+3$ & $157.2473543$ & $2.61115{\times}10^{-25}$ & $5.57763{\times}10^{-9}$ & $79.86 - 2.64039{\times}10^{-15}$ & $2.64039{\times}10^{-15}$ \\
183 & Gd nucleus & 64 & $+64$ & $157.21422$ & $2.61060{\times}10^{-25}$ & $1.18989{\times}10^{-7}$ & $79.86 - 1.23742{\times}10^{-16}$ & $1.23742{\times}10^{-16}$ \\
184 & Tb & 65 & $0$ & $158.925354$ & $2.63902{\times}10^{-25}$ & $2.63902{\times}10^{-25}$ & $0$ & $79.86$ \\
185 & Tb$^{3+}$ & 65 & $+3$ & $158.9237083$ & $2.63899{\times}10^{-25}$ & $5.57763{\times}10^{-9}$ & $79.86 - 2.66854{\times}10^{-15}$ & $2.66854{\times}10^{-15}$ \\
186 & Tb nucleus & 65 & $+65$ & $158.8900382$ & $2.63843{\times}10^{-25}$ & $1.20849{\times}10^{-7}$ & $79.86 - 1.23137{\times}10^{-16}$ & $1.23137{\times}10^{-16}$ \\
187 & Dy & 66 & $0$ & $162.5$ & $2.69838{\times}10^{-25}$ & $2.69838{\times}10^{-25}$ & $0$ & $79.86$ \\
188 & Dy$^{3+}$ & 66 & $+3$ & $162.4983543$ & $2.69835{\times}10^{-25}$ & $5.57763{\times}10^{-9}$ & $79.86 - 2.72856{\times}10^{-15}$ & $2.72856{\times}10^{-15}$ \\
189 & Dy nucleus & 66 & $+66$ & $162.4641488$ & $2.69778{\times}10^{-25}$ & $1.22708{\times}10^{-7}$ & $79.86 - 1.24000{\times}10^{-16}$ & $1.24000{\times}10^{-16}$ \\
190 & Ho & 67 & $0$ & $164.930329$ & $2.73873{\times}10^{-25}$ & $2.73873{\times}10^{-25}$ & $0$ & $79.86$ \\
191 & Ho$^{3+}$ & 67 & $+3$ & $164.9286833$ & $2.73871{\times}10^{-25}$ & $5.57763{\times}10^{-9}$ & $79.86 - 2.76937{\times}10^{-15}$ & $2.76937{\times}10^{-15}$ \\
192 & Ho nucleus & 67 & $+67$ & $164.8939426$ & $2.73813{\times}10^{-25}$ & $1.24567{\times}10^{-7}$ & $79.86 - 1.23976{\times}10^{-16}$ & $1.23976{\times}10^{-16}$ \\
193 & Er & 68 & $0$ & $167.259$ & $2.77740{\times}10^{-25}$ & $2.77740{\times}10^{-25}$ & $0$ & $79.86$ \\
194 & Er$^{3+}$ & 68 & $+3$ & $167.2573543$ & $2.77737{\times}10^{-25}$ & $5.57763{\times}10^{-9}$ & $79.86 - 2.80847{\times}10^{-15}$ & $2.80847{\times}10^{-15}$ \\
195 & Er nucleus & 68 & $+68$ & $167.2220788$ & $2.77679{\times}10^{-25}$ & $1.26426{\times}10^{-7}$ & $79.86 - 1.23877{\times}10^{-16}$ & $1.23877{\times}10^{-16}$ \\
196 & Tm & 69 & $0$ & $168.934219$ & $2.80522{\times}10^{-25}$ & $2.80522{\times}10^{-25}$ & $0$ & $79.86$ \\
197 & Tm$^{3+}$ & 69 & $+3$ & $168.9325733$ & $2.80519{\times}10^{-25}$ & $5.57763{\times}10^{-9}$ & $79.86 - 2.83660{\times}10^{-15}$ & $2.83660{\times}10^{-15}$ \\
198 & Tm nucleus & 69 & $+69$ & $168.8967632$ & $2.80460{\times}10^{-25}$ & $1.28285{\times}10^{-7}$ & $79.86 - 1.23304{\times}10^{-16}$ & $1.23304{\times}10^{-16}$ \\
199 & Yb & 70 & $0$ & $173.045$ & $2.87348{\times}10^{-25}$ & $2.87348{\times}10^{-25}$ & $0$ & $79.86$ \\
200 & Yb$^{3+}$ & 70 & $+3$ & $173.0433543$ & $2.87345{\times}10^{-25}$ & $5.57763{\times}10^{-9}$ & $79.86 - 2.90563{\times}10^{-15}$ & $2.90563{\times}10^{-15}$ \\
201 & Yb nucleus & 70 & $+70$ & $173.00701$ & $2.87285{\times}10^{-25}$ & $1.30145{\times}10^{-7}$ & $79.86 - 1.24501{\times}10^{-16}$ & $1.24501{\times}10^{-16}$ \\
202 & Lu & 71 & $0$ & $174.96669$ & $2.90539{\times}10^{-25}$ & $2.90539{\times}10^{-25}$ & $0$ & $79.86$ \\
203 & Lu$^{3+}$ & 71 & $+3$ & $174.9650443$ & $2.90536{\times}10^{-25}$ & $5.57763{\times}10^{-9}$ & $79.86 - 2.93790{\times}10^{-15}$ & $2.93790{\times}10^{-15}$ \\
204 & Lu nucleus & 71 & $+71$ & $174.9281662$ & $2.90475{\times}10^{-25}$ & $1.32004{\times}10^{-7}$ & $79.86 - 1.24110{\times}10^{-16}$ & $1.24110{\times}10^{-16}$ \\
205 & Hf & 72 & $0$ & $178.486$ & $2.96383{\times}10^{-25}$ & $2.96383{\times}10^{-25}$ & $0$ & $79.86$ \\
206 & Hf nucleus & 72 & $+72$ & $178.4469426$ & $2.96318{\times}10^{-25}$ & $1.33863{\times}10^{-7}$ & $79.86 - 1.24848{\times}10^{-16}$ & $1.24848{\times}10^{-16}$ \\
207 & Ta & 73 & $0$ & $180.94788$ & $3.00471{\times}10^{-25}$ & $3.00471{\times}10^{-25}$ & $0$ & $79.86$ \\
208 & Ta nucleus & 73 & $+73$ & $180.9082894$ & $3.00405{\times}10^{-25}$ & $1.35722{\times}10^{-7}$ & $79.86 - 1.24837{\times}10^{-16}$ & $1.24837{\times}10^{-16}$ \\
209 & W & 74 & $0$ & $183.84$ & $3.05274{\times}10^{-25}$ & $3.05274{\times}10^{-25}$ & $0$ & $79.86$ \\
210 & W nucleus & 74 & $+74$ & $183.7998766$ & $3.05207{\times}10^{-25}$ & $1.37581{\times}10^{-7}$ & $79.86 - 1.25118{\times}10^{-16}$ & $1.25118{\times}10^{-16}$ \\
211 & Re & 75 & $0$ & $186.207$ & $3.09204{\times}10^{-25}$ & $3.09204{\times}10^{-25}$ & $0$ & $79.86$ \\
212 & Re nucleus & 75 & $+75$ & $186.1663441$ & $3.09136{\times}10^{-25}$ & $1.39441{\times}10^{-7}$ & $79.86 - 1.25039{\times}10^{-16}$ & $1.25039{\times}10^{-16}$ \\
213 & Os & 76 & $0$ & $190.23$ & $3.15884{\times}10^{-25}$ & $3.15884{\times}10^{-25}$ & $0$ & $79.86$ \\
214 & Os nucleus & 76 & $+76$ & $190.1888119$ & $3.15816{\times}10^{-25}$ & $1.41300{\times}10^{-7}$ & $79.86 - 1.26060{\times}10^{-16}$ & $1.26060{\times}10^{-16}$ \\
215 & Ir & 77 & $0$ & $192.217$ & $3.19184{\times}10^{-25}$ & $3.19184{\times}10^{-25}$ & $0$ & $79.86$ \\
216 & Ir nucleus & 77 & $+77$ & $192.1752801$ & $3.19115{\times}10^{-25}$ & $1.43159{\times}10^{-7}$ & $79.86 - 1.25723{\times}10^{-16}$ & $1.25723{\times}10^{-16}$ \\
\bottomrule
\end{tabular}}
\endgroup
\clearpage
\begingroup\scriptsize\setlength{\extrarowheight}{0pt}\setlength{\tabcolsep}{3.2pt}\renewcommand{\arraystretch}{1.02}
\noindent\makebox[\textwidth][c]{%
\begin{tabular}{@{}r l r r r r r l l@{}}
\toprule
& \tableheadline{body} & $Z$ & $z$ & $m$ (u) & $m$ (kg) & $r$ (kg) & $\alpha$ (degree) & $\phi$ (degree) \\
\midrule
217 & Pt & 78 & $0$ & $195.084$ & $3.23945{\times}10^{-25}$ & $3.23945{\times}10^{-25}$ & $0$ & $79.86$ \\
218 & Pt nucleus & 78 & $+78$ & $195.0417487$ & $3.23874{\times}10^{-25}$ & $1.45018{\times}10^{-7}$ & $79.86 - 1.25962{\times}10^{-16}$ & $1.25962{\times}10^{-16}$ \\
219 & Au & 79 & $0$ & $196.96657$ & $3.27071{\times}10^{-25}$ & $3.27071{\times}10^{-25}$ & $0$ & $79.86$ \\
220 & Au nucleus & 79 & $+79$ & $196.9237876$ & $3.27000{\times}10^{-25}$ & $1.46878{\times}10^{-7}$ & $79.86 - 1.25568{\times}10^{-16}$ & $1.25568{\times}10^{-16}$ \\
221 & Hg & 80 & $0$ & $200.592$ & $3.33091{\times}10^{-25}$ & $3.33091{\times}10^{-25}$ & $0$ & $79.86$ \\
222 & Hg$^{2+}$ & 80 & $+2$ & $200.5909028$ & $3.33089{\times}10^{-25}$ & $3.71842{\times}10^{-9}$ & $79.86 - 5.05228{\times}10^{-15}$ & $5.05228{\times}10^{-15}$ \\
223 & Hg nucleus & 80 & $+80$ & $200.5486869$ & $3.33019{\times}10^{-25}$ & $1.48737{\times}10^{-7}$ & $79.86 - 1.26281{\times}10^{-16}$ & $1.26281{\times}10^{-16}$ \\
224 & Tl & 81 & $0$ & $204.38$ & $3.39381{\times}10^{-25}$ & $3.39381{\times}10^{-25}$ & $0$ & $79.86$ \\
225 & Tl$^{+}$ & 81 & $+1$ & $204.3794514$ & $3.39380{\times}10^{-25}$ & $1.85921{\times}10^{-9}$ & $79.86 - 1.02954{\times}10^{-14}$ & $1.02954{\times}10^{-14}$ \\
226 & Tl$^{3+}$ & 81 & $+3$ & $204.3783543$ & $3.39378{\times}10^{-25}$ & $5.57763{\times}10^{-9}$ & $79.86 - 3.43179{\times}10^{-15}$ & $3.43179{\times}10^{-15}$ \\
227 & Tl nucleus & 81 & $+81$ & $204.3361565$ & $3.39308{\times}10^{-25}$ & $1.50596{\times}10^{-7}$ & $79.86 - 1.27077{\times}10^{-16}$ & $1.27077{\times}10^{-16}$ \\
228 & Pb & 82 & $0$ & $207.2$ & $3.44064{\times}10^{-25}$ & $3.44064{\times}10^{-25}$ & $0$ & $79.86$ \\
229 & Pb$^{2+}$ & 82 & $+2$ & $207.1989028$ & $3.44062{\times}10^{-25}$ & $3.71842{\times}10^{-9}$ & $79.86 - 5.21872{\times}10^{-15}$ & $5.21872{\times}10^{-15}$ \\
230 & Pb$^{4+}$ & 82 & $+4$ & $207.1978057$ & $3.44060{\times}10^{-25}$ & $7.43684{\times}10^{-9}$ & $79.86 - 2.60935{\times}10^{-15}$ & $2.60935{\times}10^{-15}$ \\
231 & Pb nucleus & 82 & $+82$ & $207.1556266$ & $3.43990{\times}10^{-25}$ & $1.52455{\times}10^{-7}$ & $79.86 - 1.27259{\times}10^{-16}$ & $1.27259{\times}10^{-16}$ \\
232 & Bi & 83 & $0$ & $208.9804$ & $3.47020{\times}10^{-25}$ & $3.47020{\times}10^{-25}$ & $0$ & $79.86$ \\
233 & Bi nucleus & 83 & $+83$ & $208.935497$ & $3.46946{\times}10^{-25}$ & $1.54314{\times}10^{-7}$ & $79.86 - 1.26806{\times}10^{-16}$ & $1.26806{\times}10^{-16}$ \\
234 & Po & 84 & $0$ & $209.0$ & $3.47053{\times}10^{-25}$ & $3.47053{\times}10^{-25}$ & $0$ & $79.86$ \\
235 & Po nucleus & 84 & $+84$ & $208.9545679$ & $3.46977{\times}10^{-25}$ & $1.56174{\times}10^{-7}$ & $79.86 - 1.25308{\times}10^{-16}$ & $1.25308{\times}10^{-16}$ \\
236 & At & 85 & $0$ & $210.0$ & $3.48713{\times}10^{-25}$ & $3.48713{\times}10^{-25}$ & $0$ & $79.86$ \\
237 & At nucleus & 85 & $+85$ & $209.9540391$ & $3.48637{\times}10^{-25}$ & $1.58033{\times}10^{-7}$ & $79.86 - 1.24426{\times}10^{-16}$ & $1.24426{\times}10^{-16}$ \\
238 & Rn & 86 & $0$ & $222.0$ & $3.68640{\times}10^{-25}$ & $3.68640{\times}10^{-25}$ & $0$ & $79.86$ \\
239 & Rn nucleus & 86 & $+86$ & $221.9535108$ & $3.68562{\times}10^{-25}$ & $1.59892{\times}10^{-7}$ & $79.86 - 1.30008{\times}10^{-16}$ & $1.30008{\times}10^{-16}$ \\
240 & Fr & 87 & $0$ & $223.0$ & $3.70300{\times}10^{-25}$ & $3.70300{\times}10^{-25}$ & $0$ & $79.86$ \\
241 & Fr$^{+}$ & 87 & $+1$ & $222.9994514$ & $3.70299{\times}10^{-25}$ & $1.85921{\times}10^{-9}$ & $79.86 - 1.12334{\times}10^{-14}$ & $1.12334{\times}10^{-14}$ \\
242 & Fr nucleus & 87 & $+87$ & $222.9529828$ & $3.70222{\times}10^{-25}$ & $1.61751{\times}10^{-7}$ & $79.86 - 1.29092{\times}10^{-16}$ & $1.29092{\times}10^{-16}$ \\
243 & Ra & 88 & $0$ & $226.0$ & $3.75282{\times}10^{-25}$ & $3.75282{\times}10^{-25}$ & $0$ & $79.86$ \\
244 & Ra$^{2+}$ & 88 & $+2$ & $225.9989028$ & $3.75280{\times}10^{-25}$ & $3.71842{\times}10^{-9}$ & $79.86 - 5.69224{\times}10^{-15}$ & $5.69224{\times}10^{-15}$ \\
245 & Ra nucleus & 88 & $+88$ & $225.9524553$ & $3.75203{\times}10^{-25}$ & $1.63610{\times}10^{-7}$ & $79.86 - 1.29342{\times}10^{-16}$ & $1.29342{\times}10^{-16}$ \\
246 & Ac & 89 & $0$ & $227.0$ & $3.76942{\times}10^{-25}$ & $3.76942{\times}10^{-25}$ & $0$ & $79.86$ \\
247 & Ac$^{3+}$ & 89 & $+3$ & $226.9983543$ & $3.76940{\times}10^{-25}$ & $5.57763{\times}10^{-9}$ & $79.86 - 3.81161{\times}10^{-15}$ & $3.81161{\times}10^{-15}$ \\
248 & Ac nucleus & 89 & $+89$ & $226.9519282$ & $3.76863{\times}10^{-25}$ & $1.65470{\times}10^{-7}$ & $79.86 - 1.28455{\times}10^{-16}$ & $1.28455{\times}10^{-16}$ \\
249 & Th & 90 & $0$ & $232.0377$ & $3.85308{\times}10^{-25}$ & $3.85308{\times}10^{-25}$ & $0$ & $79.86$ \\
250 & Th nucleus & 90 & $+90$ & $231.9891015$ & $3.85227{\times}10^{-25}$ & $1.67329{\times}10^{-7}$ & $79.86 - 1.29847{\times}10^{-16}$ & $1.29847{\times}10^{-16}$ \\
251 & Pa & 91 & $0$ & $231.03588$ & $3.83644{\times}10^{-25}$ & $3.83644{\times}10^{-25}$ & $0$ & $79.86$ \\
252 & Pa nucleus & 91 & $+91$ & $230.9867552$ & $3.83563{\times}10^{-25}$ & $1.69188{\times}10^{-7}$ & $79.86 - 1.27865{\times}10^{-16}$ & $1.27865{\times}10^{-16}$ \\
253 & U & 92 & $0$ & $238.02891$ & $3.95256{\times}10^{-25}$ & $3.95256{\times}10^{-25}$ & $0$ & $79.86$ \\
254 & U nucleus & 92 & $+92$ & $237.9792594$ & $3.95174{\times}10^{-25}$ & $1.71047{\times}10^{-7}$ & $79.86 - 1.30304{\times}10^{-16}$ & $1.30304{\times}10^{-16}$ \\
255 & Np & 93 & $0$ & $237.0$ & $3.93548{\times}10^{-25}$ & $3.93548{\times}10^{-25}$ & $0$ & $79.86$ \\
256 & Np nucleus & 93 & $+93$ & $236.949824$ & $3.93464{\times}10^{-25}$ & $1.72906{\times}10^{-7}$ & $79.86 - 1.28345{\times}10^{-16}$ & $1.28345{\times}10^{-16}$ \\
257 & Pu & 94 & $0$ & $244.0$ & $4.05172{\times}10^{-25}$ & $4.05172{\times}10^{-25}$ & $0$ & $79.86$ \\
258 & Pu nucleus & 94 & $+94$ & $243.9492991$ & $4.05087{\times}10^{-25}$ & $1.74766{\times}10^{-7}$ & $79.86 - 1.30731{\times}10^{-16}$ & $1.30731{\times}10^{-16}$ \\
259 & Am & 95 & $0$ & $243.0$ & $4.03511{\times}10^{-25}$ & $4.03511{\times}10^{-25}$ & $0$ & $79.86$ \\
260 & Am nucleus & 95 & $+95$ & $242.9487747$ & $4.03426{\times}10^{-25}$ & $1.76625{\times}10^{-7}$ & $79.86 - 1.28824{\times}10^{-16}$ & $1.28824{\times}10^{-16}$ \\
261 & Cm & 96 & $0$ & $247.0$ & $4.10153{\times}10^{-25}$ & $4.10153{\times}10^{-25}$ & $0$ & $79.86$ \\
262 & Cm nucleus & 96 & $+96$ & $246.9482507$ & $4.10067{\times}10^{-25}$ & $1.78484{\times}10^{-7}$ & $79.86 - 1.29581{\times}10^{-16}$ & $1.29581{\times}10^{-16}$ \\
263 & Bk & 97 & $0$ & $247.0$ & $4.10153{\times}10^{-25}$ & $4.10153{\times}10^{-25}$ & $0$ & $79.86$ \\
264 & Bk nucleus & 97 & $+97$ & $246.9477272$ & $4.10066{\times}10^{-25}$ & $1.80343{\times}10^{-7}$ & $79.86 - 1.28245{\times}10^{-16}$ & $1.28245{\times}10^{-16}$ \\
265 & Cf & 98 & $0$ & $251.0$ & $4.16795{\times}10^{-25}$ & $4.16795{\times}10^{-25}$ & $0$ & $79.86$ \\
266 & Cf nucleus & 98 & $+98$ & $250.9472042$ & $4.16708{\times}10^{-25}$ & $1.82202{\times}10^{-7}$ & $79.86 - 1.28992{\times}10^{-16}$ & $1.28992{\times}10^{-16}$ \\
267 & Es & 99 & $0$ & $252.0$ & $4.18456{\times}10^{-25}$ & $4.18456{\times}10^{-25}$ & $0$ & $79.86$ \\
268 & Es nucleus & 99 & $+99$ & $251.9466817$ & $4.18367{\times}10^{-25}$ & $1.84062{\times}10^{-7}$ & $79.86 - 1.28198{\times}10^{-16}$ & $1.28198{\times}10^{-16}$ \\
269 & Fm & 100 & $0$ & $257.0$ & $4.26759{\times}10^{-25}$ & $4.26759{\times}10^{-25}$ & $0$ & $79.86$ \\
270 & Fm nucleus & 100 & $+100$ & $256.9461596$ & $4.26669{\times}10^{-25}$ & $1.85921{\times}10^{-7}$ & $79.86 - 1.29434{\times}10^{-16}$ & $1.29434{\times}10^{-16}$ \\
271 & Md & 101 & $0$ & $258.0$ & $4.28419{\times}10^{-25}$ & $4.28419{\times}10^{-25}$ & $0$ & $79.86$ \\
272 & Md nucleus & 101 & $+101$ & $257.9456381$ & $4.28329{\times}10^{-25}$ & $1.87780{\times}10^{-7}$ & $79.86 - 1.28651{\times}10^{-16}$ & $1.28651{\times}10^{-16}$ \\
273 & No & 102 & $0$ & $259.0$ & $4.30080{\times}10^{-25}$ & $4.30080{\times}10^{-25}$ & $0$ & $79.86$ \\
274 & No nucleus & 102 & $+102$ & $258.9451171$ & $4.29988{\times}10^{-25}$ & $1.89639{\times}10^{-7}$ & $79.86 - 1.27883{\times}10^{-16}$ & $1.27883{\times}10^{-16}$ \\
275 & Lr & 103 & $0$ & $266.0$ & $4.41703{\times}10^{-25}$ & $4.41703{\times}10^{-25}$ & $0$ & $79.86$ \\
276 & Lr nucleus & 103 & $+103$ & $265.9445966$ & $4.41611{\times}10^{-25}$ & $1.91499{\times}10^{-7}$ & $79.86 - 1.30065{\times}10^{-16}$ & $1.30065{\times}10^{-16}$ \\
\bottomrule
\end{tabular}}
\endgroup
\clearpage
\begingroup\scriptsize\setlength{\extrarowheight}{0pt}\setlength{\tabcolsep}{3.2pt}\renewcommand{\arraystretch}{1.02}
\noindent\makebox[\textwidth][c]{%
\begin{tabular}{@{}r l r r r r r l l@{}}
\toprule
& \tableheadline{body} & $Z$ & $z$ & $m$ (u) & $m$ (kg) & $r$ (kg) & $\alpha$ (degree) & $\phi$ (degree) \\
\midrule
277 & Rf & 104 & $0$ & $267.0$ & $4.43364{\times}10^{-25}$ & $4.43364{\times}10^{-25}$ & $0$ & $79.86$ \\
278 & Rf nucleus & 104 & $+104$ & $266.9440766$ & $4.43271{\times}10^{-25}$ & $1.93358{\times}10^{-7}$ & $79.86 - 1.29298{\times}10^{-16}$ & $1.29298{\times}10^{-16}$ \\
279 & Db & 105 & $0$ & $268.0$ & $4.45024{\times}10^{-25}$ & $4.45024{\times}10^{-25}$ & $0$ & $79.86$ \\
280 & Db nucleus & 105 & $+105$ & $267.9435572$ & $4.44931{\times}10^{-25}$ & $1.95217{\times}10^{-7}$ & $79.86 - 1.28547{\times}10^{-16}$ & $1.28547{\times}10^{-16}$ \\
281 & Sg & 106 & $0$ & $267.0$ & $4.43364{\times}10^{-25}$ & $4.43364{\times}10^{-25}$ & $0$ & $79.86$ \\
282 & Sg nucleus & 106 & $+106$ & $266.9430383$ & $4.43269{\times}10^{-25}$ & $1.97076{\times}10^{-7}$ & $79.86 - 1.26858{\times}10^{-16}$ & $1.26858{\times}10^{-16}$ \\
283 & Bh & 107 & $0$ & $270.0$ & $4.48346{\times}10^{-25}$ & $4.48346{\times}10^{-25}$ & $0$ & $79.86$ \\
284 & Bh nucleus & 107 & $+107$ & $269.94252$ & $4.48250{\times}10^{-25}$ & $1.98935{\times}10^{-7}$ & $79.86 - 1.27085{\times}10^{-16}$ & $1.27085{\times}10^{-16}$ \\
285 & Hs & 108 & $0$ & $271.0$ & $4.50006{\times}10^{-25}$ & $4.50006{\times}10^{-25}$ & $0$ & $79.86$ \\
286 & Hs nucleus & 108 & $+108$ & $270.9420022$ & $4.49910{\times}10^{-25}$ & $2.00795{\times}10^{-7}$ & $79.86 - 1.26374{\times}10^{-16}$ & $1.26374{\times}10^{-16}$ \\
287 & Mt & 109 & $0$ & $278.0$ & $4.61630{\times}10^{-25}$ & $4.61630{\times}10^{-25}$ & $0$ & $79.86$ \\
288 & Mt nucleus & 109 & $+109$ & $277.941485$ & $4.61533{\times}10^{-25}$ & $2.02654{\times}10^{-7}$ & $79.86 - 1.28450{\times}10^{-16}$ & $1.28450{\times}10^{-16}$ \\
289 & Ds & 110 & $0$ & $281.0$ & $4.66611{\times}10^{-25}$ & $4.66611{\times}10^{-25}$ & $0$ & $79.86$ \\
290 & Ds nucleus & 110 & $+110$ & $280.9409684$ & $4.66513{\times}10^{-25}$ & $2.04513{\times}10^{-7}$ & $79.86 - 1.28656{\times}10^{-16}$ & $1.28656{\times}10^{-16}$ \\
291 & Rg & 111 & $0$ & $282.0$ & $4.68272{\times}10^{-25}$ & $4.68272{\times}10^{-25}$ & $0$ & $79.86$ \\
292 & Rg nucleus & 111 & $+111$ & $281.9404524$ & $4.68173{\times}10^{-25}$ & $2.06372{\times}10^{-7}$ & $79.86 - 1.27950{\times}10^{-16}$ & $1.27950{\times}10^{-16}$ \\
293 & Cn & 112 & $0$ & $285.0$ & $4.73254{\times}10^{-25}$ & $4.73254{\times}10^{-25}$ & $0$ & $79.86$ \\
294 & Cn nucleus & 112 & $+112$ & $284.9399369$ & $4.73154{\times}10^{-25}$ & $2.08231{\times}10^{-7}$ & $79.86 - 1.28157{\times}10^{-16}$ & $1.28157{\times}10^{-16}$ \\
295 & Nh & 113 & $0$ & $286.0$ & $4.74914{\times}10^{-25}$ & $4.74914{\times}10^{-25}$ & $0$ & $79.86$ \\
296 & Nh nucleus & 113 & $+113$ & $285.9394221$ & $4.74814{\times}10^{-25}$ & $2.10091{\times}10^{-7}$ & $79.86 - 1.27468{\times}10^{-16}$ & $1.27468{\times}10^{-16}$ \\
297 & Fl & 114 & $0$ & $289.0$ & $4.79896{\times}10^{-25}$ & $4.79896{\times}10^{-25}$ & $0$ & $79.86$ \\
298 & Fl nucleus & 114 & $+114$ & $288.9389079$ & $4.79794{\times}10^{-25}$ & $2.11950{\times}10^{-7}$ & $79.86 - 1.27676{\times}10^{-16}$ & $1.27676{\times}10^{-16}$ \\
299 & Mc & 115 & $0$ & $290.0$ & $4.81556{\times}10^{-25}$ & $4.81556{\times}10^{-25}$ & $0$ & $79.86$ \\
300 & Mc nucleus & 115 & $+115$ & $289.9383944$ & $4.81454{\times}10^{-25}$ & $2.13809{\times}10^{-7}$ & $79.86 - 1.27003{\times}10^{-16}$ & $1.27003{\times}10^{-16}$ \\
301 & Lv & 116 & $0$ & $293.0$ & $4.86538{\times}10^{-25}$ & $4.86538{\times}10^{-25}$ & $0$ & $79.86$ \\
302 & Lv nucleus & 116 & $+116$ & $292.9378814$ & $4.86435{\times}10^{-25}$ & $2.15668{\times}10^{-7}$ & $79.86 - 1.27211{\times}10^{-16}$ & $1.27211{\times}10^{-16}$ \\
303 & Ts & 117 & $0$ & $294.0$ & $4.88198{\times}10^{-25}$ & $4.88198{\times}10^{-25}$ & $0$ & $79.86$ \\
304 & Ts nucleus & 117 & $+117$ & $293.9373691$ & $4.88094{\times}10^{-25}$ & $2.17527{\times}10^{-7}$ & $79.86 - 1.26554{\times}10^{-16}$ & $1.26554{\times}10^{-16}$ \\
305 & Og & 118 & $0$ & $294.0$ & $4.88198{\times}10^{-25}$ & $4.88198{\times}10^{-25}$ & $0$ & $79.86$ \\
306 & Og nucleus & 118 & $+118$ & $293.9368575$ & $4.88094{\times}10^{-25}$ & $2.19387{\times}10^{-7}$ & $79.86 - 1.25481{\times}10^{-16}$ & $1.25481{\times}10^{-16}$ \\
\bottomrule
\end{tabular}}
\endgroup
\clearpage

.
% Every law below was checked with verify_general_number_algebra.py and by an independent review.
\providecommand{\Jp}{\mathcal{J}_{\bar{\phi}}}                      % the Turning Unit
\providecommand{\Ht}{\operatorname{Ht}}                              % the height above the real axis
\providecommand{\Sh}{\operatorname{Sh}}                              % the shadow on the real axis
\providecommand{\Com}{\operatorname{Com}}                            % the Composer part (Definition 4.11)
% start a new page if fewer than #1 lines are left (the same command as in Section 4.9)
\makeatletter
\providecommand{\gnkeep}[1]{\par\penalty\z@\begingroup\dimen@=\pagegoal\advance\dimen@-\pagetotal
  \ifdim\pagetotal>\z@\ifdim\dimen@<#1\baselineskip\newpage\fi\fi\endgroup}
\makeatother
% one row of a table of laws: the name of the law, then the law and its complex case side by side
\providecommand{\gnlaw}[3]{\multicolumn{2}{@{}l@{}}{\small\itshape #1}\\* #2 & #3 \\ \addlinespace[4pt]}
% the two columns of the tables
\newenvironment{gnlaws}{%
  \begin{longtable}{@{}>{\small\raggedright\arraybackslash}p{0.62\textwidth}@{\hspace{0.04\textwidth}}>{\small\raggedright\arraybackslash}p{0.32\textwidth}@{}}
  \toprule
  \multicolumn{1}{@{}l}{\spacedlowsmallcaps{in a relative plane}} & \multicolumn{1}{l@{}}{\spacedlowsmallcaps{the complex case}} \\
  \midrule
  \endhead
  \bottomrule
  \endlastfoot}
  {\end{longtable}}
% the same two columns for the last section, with a wider left column for its longer laws;
% there a line breaks only between two formulas, never inside one
\newenvironment{gnlawswide}{%
  \begin{longtable}{@{}>{\small\raggedright\arraybackslash\relpenalty=10000 \binoppenalty=10000 }p{0.64\textwidth}@{\hspace{0.03\textwidth}}>{\small\raggedright\arraybackslash}p{0.32\textwidth}@{}}
  \toprule
  \multicolumn{1}{@{}l}{\spacedlowsmallcaps{in a relative plane}} & \multicolumn{1}{l@{}}{\spacedlowsmallcaps{the complex case}} \\
  \midrule
  \endhead
  \bottomrule
  \endlastfoot}
  {\end{longtable}}

This appendix lists the laws of \Cref{sec:gn-algebra}, each beside its complex case, in the order of Appendix~\ref{app:complex}. The letters follow that appendix, so that each line can be read against it. The same laws are collected in the formulary of \citet[Appendix~A]{DerrazGNPreprint2026}.
\begin{itemize}
  \item $z = a + b\Gp$ and $w = c + d\Gp$ are General Numbers of one true relative plane $\GNp$ (Remark~\ref{rem:line-no-plane}).
  \item $|\Gp| = |\cos\pb + \sin\pb|$ is the length of the Composer Unit, $\up = \Gp/|\Gp|$ is its direction, and $\Jp$ is the Turning Unit (Definition~\ref{def:gn-turning-unit}).
  \item $\theta$ is the angle of a number, the studied angle of Definition~\ref{def:general-polar}, and $\pb - \theta$ is its projection angle. For short, $C_\theta = \cos(\theta, \pb - \theta)$ and $S_\theta = \sin(\theta, \pb - \theta)$.
\end{itemize}
The right-hand column, the complex case, gives the same law at $\pb = \pi/2$. There $|\Gp| = 1$, the three numbers $\Gp$, $\up$ and $\Jp$ are all $i$, and $C_\theta$ and $S_\theta$ are $\cos\theta$ and $\sin\theta$ (Corollary~\ref{cor:classical}). Otherwise the two columns differ only in names, which the right-hand column writes as Appendix~\ref{app:complex} does: $\operatorname{Re}$ and $\operatorname{Im}$ for the parts, and $r$ for $|z|$ in the polar forms. A few laws are not stated in \Cref{sec:gn-algebra}; a line under each table says where they come from.

% keep the heading with its table: break the page here if fewer than 12 lines are left
\gnkeep{12}
%------------------------------------------------
\section{The two units}\label{app:f-units}
%------------------------------------------------

\begin{gnlaws}
\gnlaw{the Rule of the Square (Proposition~\ref{prop:square-rule})}
  {$\Gp^2 = 2|\Gp|\cos\pb\;\Gp - |\Gp|^2$, \quad $\up^2 = 2\cos\pb\;\up - 1$}
  {$i^2 = -1$}
\gnlaw{the Composer Unit and its conjugate (\Cref{thm:gn-conjugate})}
  {$\Gp + \overline{\Gp} = 2|\Gp|\cos\pb$, \quad $\Gp\,\overline{\Gp} = |\Gp|^2$}
  {$i + (-i) = 0$, \quad $i\,(-i) = 1$}
\gnlaw{the Turning Unit (\Cref{thm:gn-turning-unit})}
  {$\Jp = \dfrac{\up - \cos\pb}{\sin\pb} = -\cot\pb + \dfrac{\Gp}{|\Gp|\sin\pb}$, \quad $\Jp^2 = -1$}
  {$i^2 = -1$}
\gnlaw{the square roots of $-1$}
  {$\Jp$ and $-\Jp$}
  {$i$ and $-i$}
\gnlaw{shadow and height (Proposition~\ref{prop:gn-shadow-height})}
  {$\Sh(z) = a + b|\Gp|\cos\pb$, \quad $\Ht(z) = b|\Gp|\sin\pb$}
  {$a$ and $b$, the real and imaginary parts}
\gnlaw{the two frames}
  {$z = a + b\Gp = \Sh(z) + \Ht(z)\,\Jp$}
  {$z = a + bi$}
\gnlaw{back to the leaning frame}
  {$a = \Sh(z) - \Ht(z)\cot\pb$, \quad $b = \dfrac{\Ht(z)}{|\Gp|\sin\pb}$}
  {$a$ and $b$}
\gnlaw{from the complex plane}
  {$x + iy = (x - y\cot\pb) + \dfrac{y}{|\Gp|\sin\pb}\,\Gp$}
  {$x + iy$}
\gnlaw{to another true relative plane $\mathbb{G}_{\bar{\psi}}$}
  {$z = a' + b'\mathcal{G}_{\bar{\psi}}$ with $b' = \dfrac{\Ht(z)}{|\mathcal{G}_{\bar{\psi}}|\sin\bar{\psi}}$, \quad $a' = \Sh(z) - \Ht(z)\cot\bar{\psi}$}
  {at $\bar{\psi} = \pi/2$: \newline $a' = \Sh(z)$, \quad $b' = \Ht(z)$}
\end{gnlaws}
The shadow and the height are read on the page, so they do not depend on the frame: a number keeps them when it is written in the complex plane or in another relative plane, and the last two lines follow from Proposition~\ref{prop:gn-shadow-height}.

% keep the heading with its table: break the page here if fewer than 12 lines are left
\gnkeep{12}
%------------------------------------------------
\section{Arithmetic}\label{app:f-arithmetic}
%------------------------------------------------

\begin{gnlaws}
\gnlaw{sum}
  {$(a + b\Gp) + (c + d\Gp) = (a + c) + (b + d)\,\Gp$}
  {$(a + c) + (b + d)\,i$}
\gnlaw{product (Proposition~\ref{prop:un-real-product})}
  {$(a + b\Gp)(c + d\Gp)$ \newline $\quad = (ac - |\Gp|^2 bd) + (ad + bc + 2|\Gp|\cos\pb\;bd)\,\Gp$}
  {$(ac - bd) + (ad + bc)\,i$}
\gnlaw{the rules of arithmetic (\Cref{thm:gn-field})}
  {order, grouping, distribution; $0$, $1$, opposites and inverses: a field}
  {the same}
\gnlaw{conjugate (Definition~\ref{def:gn-conjugate})}
  {$\bar{z} = (a + 2b|\Gp|\cos\pb) - b\Gp$, the mirror image in the real axis}
  {$\bar{z} = a - bi$}
\gnlaw{the parts of the conjugate (Definitions~\ref{def:gn-parts} and~\ref{def:gn-conjugate})}
  {$\Re(\bar{z}) = a + 2b|\Gp|\cos\pb$, \quad $\Com(\bar{z}) = -b\Gp = -\Com(z)$}
  {$\operatorname{Re}\bar{z} = \operatorname{Re} z$, \newline $\operatorname{Im}\bar{z} = -\operatorname{Im} z$}
\gnlaw{a number times its conjugate}
  {$z\bar{z} = a^2 + 2ab|\Gp|\cos\pb + b^2|\Gp|^2 = |z|^2$}
  {$z\bar{z} = a^2 + b^2$}
\gnlaw{a number plus or minus its conjugate}
  {$z + \bar{z} = 2\Sh(z)$, \quad $z - \bar{z} = 2\Ht(z)\,\Jp$}
  {$z + \bar{z} = 2a$, \quad $z - \bar{z} = 2bi$}
\gnlaw{the conjugate of a sum and of a product}
  {$\overline{z + w} = \bar{z} + \bar{w}$, \quad $\overline{zw} = \bar{z}\,\bar{w}$, \quad $\bar{\bar{z}} = z$}
  {the same}
\gnlaw{inverse (\Cref{thm:gn-division})}
  {$\dfrac{1}{z} = \dfrac{(a + 2b|\Gp|\cos\pb) - b\Gp}{a^2 + 2ab|\Gp|\cos\pb + b^2|\Gp|^2}$, \quad $\dfrac{1}{\Gp} = \dfrac{2|\Gp|\cos\pb - \Gp}{|\Gp|^2}$}
  {$\dfrac{1}{z} = \dfrac{a - bi}{a^2 + b^2}$, \quad $\dfrac{1}{i} = -i$}
\gnlaw{quotient (\Cref{thm:gn-division})}
  {$\dfrac{a + b\Gp}{c + d\Gp} = \dfrac{(ac + 2ad|\Gp|\cos\pb + bd|\Gp|^2) + (bc - ad)\,\Gp}{c^2 + 2cd|\Gp|\cos\pb + d^2|\Gp|^2}$}
  {$\dfrac{(ac + bd) + (bc - ad)\,i}{c^2 + d^2}$}
\gnlaw{the skew flip, which is not the conjugate (Remark~\ref{rem:gn-skew-flip})}
  {$(a + b\Gp)(a - b\Gp) = a^2 + b^2|\Gp|^2 - 2b^2|\Gp|\cos\pb\;\Gp$}
  {$(a + bi)(a - bi) = a^2 + b^2$}
\end{gnlaws}

% keep the heading with its table: break the page here if fewer than 12 lines are left
\gnkeep{12}
%------------------------------------------------
\section{The plane}\label{app:f-plane}
%------------------------------------------------

\begin{gnlaws}
\gnlaw{magnitude (\Cref{thm:triangle})}
  {$|z| = \sqrt{a^2 + 2ab|\Gp|\cos\pb + b^2|\Gp|^2} = \sqrt{\Sh(z)^2 + \Ht(z)^2}$}
  {$|z| = \sqrt{a^2 + b^2}$}
\gnlaw{the magnitude of a product (\Cref{thm:gn-two-squares})}
  {$|zw| = |z|\,|w|$}
  {the same}
\gnlaw{the two-squares identity, with $zw = p + q\Gp$}
  {$(a^2 + 2ab|\Gp|\cos\pb + b^2|\Gp|^2) \times {}$ \newline $(c^2 + 2cd|\Gp|\cos\pb + d^2|\Gp|^2) = {}$ \newline $p^2 + 2pq|\Gp|\cos\pb + q^2|\Gp|^2$}
  {$(a^2 + b^2)(c^2 + d^2) = {}$ \newline $(ac - bd)^2 + (ad + bc)^2$}
\gnlaw{angle (\Cref{thm:phase})}
  {$\tan\theta = \dfrac{\Ht(z)}{\Sh(z)} = \dfrac{b|\Gp|\sin\pb}{a + b|\Gp|\cos\pb}$}
  {$\tan\theta = \dfrac{b}{a}$}
\gnlaw{polar form (\Cref{thm:gn-polar})}
  {$z = |z|\,\bigl[\cos(\theta, \pb - \theta) + \sin(\theta, \pb - \theta)\,\up\bigr] = |z|\,(C_\theta + S_\theta\,\up)$}
  {$z = r(\cos\theta + i\sin\theta)$}
\gnlaw{the parts in polar form}
  {$a = |z|\,C_\theta$, \quad $b\,|\Gp| = |z|\,S_\theta$}
  {$a = r\cos\theta$, \quad $b = r\sin\theta$}
\gnlaw{the general Pythagoras identity}
  {$C_\theta^2 + 2\cos\pb\;C_\theta S_\theta + S_\theta^2 = 1$}
  {$\cos^2\theta + \sin^2\theta = 1$}
\gnlaw{the opposite angle}
  {$C_{-\theta} = C_\theta + 2\cos\pb\;S_\theta$, \quad $S_{-\theta} = -S_\theta$}
  {$\cos(-\theta) = \cos\theta$, \newline $\sin(-\theta) = -\sin\theta$}
\gnlaw{the conjugate in polar form}
  {$\bar{z} = |z|\,(C_{-\theta} + S_{-\theta}\,\up)$: the same magnitude, the opposite angle}
  {$\bar{z} = r(\cos\theta - i\sin\theta)$}
\gnlaw{the mirror in the bisector of the two axes}
  {$z \mapsto \up\,\bar{z}$ sends $x + y\,\up$ to $y + x\,\up$}
  {$z \mapsto i\bar{z}$ sends $a + bi$ to $b + ai$}
\gnlaw{reciprocity (Appendix~\ref{app:general-sinus})}
  {$C_{\pb - \theta} = S_\theta$ and $S_{\pb - \theta} = C_\theta$: the unit number at the angle $\pb - \theta$ has the two parts of the one at $\theta$, swapped}
  {$\cos(\tfrac{\pi}{2} - \theta) = \sin\theta$, \newline $\sin(\tfrac{\pi}{2} - \theta) = \cos\theta$}
\end{gnlaws}
The angle is the one of \Cref{thm:phase}, whose ratio $\eta = |\Gp|\,b/a$ gives $\eta\sin\pb/(1 + \eta\cos\pb) = \Ht(z)/\Sh(z)$. The opposite angle follows from $\sin(\pb + \theta) = \sin(\pb - \theta) + 2\cos\pb\sin\theta$. The mirror in the line at the angle $\pb/2$, which halves the angle between the two axes, keeps the length of a number and turns its angle $\theta$ into $\pb - \theta$; reciprocity, $\sin(\alpha, \phi) = \cos(\phi, \alpha)$, says the same in the General Sinus.

% keep the heading with its table: break the page here if fewer than 12 lines are left
\gnkeep{12}
%------------------------------------------------
\section{Multiplying means turning}\label{app:f-turning}
%------------------------------------------------

\begin{gnlaws}
\gnlaw{product in polar form (\Cref{thm:gn-turning})}
  {$|z|\,(C_{\theta_1} + S_{\theta_1}\up) \cdot |w|\,(C_{\theta_2} + S_{\theta_2}\up) = |z|\,|w|\,(C_{\theta_1 + \theta_2} + S_{\theta_1 + \theta_2}\up)$}
  {$r_1r_2\bigl(\cos(\theta_1 + \theta_2) + i\sin(\theta_1 + \theta_2)\bigr)$}
\gnlaw{quotient in polar form}
  {$\dfrac{z}{w} = \dfrac{|z|}{|w|}\,(C_{\theta_1 - \theta_2} + S_{\theta_1 - \theta_2}\up)$}
  {$\dfrac{r_1}{r_2}\bigl(\cos(\theta_1 - \theta_2) + i\sin(\theta_1 - \theta_2)\bigr)$}
\gnlaw{addition along the Composer axis, cosine}
  {$C_{\alpha + \beta} = C_\alpha C_\beta - S_\alpha S_\beta$}
  {$\cos(\alpha + \beta) = {}$ \newline $\cos\alpha\cos\beta - \sin\alpha\sin\beta$}
\gnlaw{addition along the Composer axis, sine}
  {$S_{\alpha + \beta} = C_\alpha S_\beta + S_\alpha C_\beta + 2\cos\pb\;S_\alpha S_\beta$}
  {$\sin(\alpha + \beta) = {}$ \newline $\sin\alpha\cos\beta + \cos\alpha\sin\beta$}
\gnlaw{subtraction, cosine}
  {$C_{\alpha - \beta} = C_\alpha C_\beta + S_\alpha S_\beta + 2\cos\pb\;C_\alpha S_\beta$}
  {$\cos(\alpha - \beta) = {}$ \newline $\cos\alpha\cos\beta + \sin\alpha\sin\beta$}
\gnlaw{subtraction, sine}
  {$S_{\alpha - \beta} = S_\alpha C_\beta - C_\alpha S_\beta$}
  {$\sin(\alpha - \beta) = {}$ \newline $\sin\alpha\cos\beta - \cos\alpha\sin\beta$}
\gnlaw{double angle}
  {$C_{2\alpha} = C_\alpha^2 - S_\alpha^2$, \quad $S_{2\alpha} = 2S_\alpha\,(C_\alpha + \cos\pb\;S_\alpha)$}
  {$\cos 2\alpha = \cos^2\alpha - \sin^2\alpha$, \newline $\sin 2\alpha = 2\sin\alpha\cos\alpha$}
\gnlaw{a turn by $\pb$, the product by $\up$}
  {$C_{\theta + \pb} = -S_\theta$, \quad $S_{\theta + \pb} = C_\theta + 2\cos\pb\;S_\theta$}
  {$\cos(\theta + \tfrac{\pi}{2}) = -\sin\theta$, \newline $\sin(\theta + \tfrac{\pi}{2}) = \cos\theta$}
\gnlaw{a quarter turn, the product by $\Jp$}
  {$C_{\theta + \pi/2} = -C_\theta\cot\pb - \dfrac{S_\theta}{\sin\pb}$, \quad $S_{\theta + \pi/2} = \dfrac{C_\theta}{\sin\pb} + S_\theta\cot\pb$}
  {$\cos(\theta + \tfrac{\pi}{2}) = -\sin\theta$, \newline $\sin(\theta + \tfrac{\pi}{2}) = \cos\theta$}
\gnlaw{a half turn, the product by $-1$}
  {$C_{\theta + \pi} = -C_\theta$, \quad $S_{\theta + \pi} = -S_\theta$}
  {$\cos(\theta + \pi) = -\cos\theta$, \newline $\sin(\theta + \pi) = -\sin\theta$}
\end{gnlaws}
The subtraction formulas are the addition formulas with the opposite angle of \Cref{app:f-plane}, and the double angle is $\beta = \alpha$. A turn by $\pb$ is the product by $\up = e^{\pb\Jp}$, a quarter turn the product by $\Jp$, and a half turn the product by $e^{\pi\Jp} = -1$; the last is also the $\pi$-addition identity of Appendix~\ref{app:general-sinus}. In the complex plane the turn by $\pb$ and the quarter turn are one and the same.

% keep the heading with its table: break the page here if fewer than 12 lines are left
\gnkeep{12}
%------------------------------------------------
\section{Euler's formula and de Moivre's formula}\label{app:f-euler}
%------------------------------------------------

\begin{gnlaws}
\gnlaw{Euler's formula (\Cref{thm:gn-euler})}
  {$e^{\theta\Jp} = \cos(\theta, \pb - \theta) + \sin(\theta, \pb - \theta)\,\up = C_\theta + S_\theta\,\up$}
  {$e^{i\theta} = \cos\theta + i\sin\theta$}
\gnlaw{exponential form}
  {$z = |z|\,e^{\theta\Jp}$, \quad $\bar{z} = |z|\,e^{-\theta\Jp}$}
  {$z = re^{i\theta}$, \quad $\bar{z} = re^{-i\theta}$}
\gnlaw{the two units as exponentials}
  {$\Gp = |\Gp|\,e^{\pb\Jp}$, \quad $\Jp = e^{(\pi/2)\Jp}$}
  {$i = e^{i\pi/2}$}
\gnlaw{Euler's identity}
  {$e^{\pi\Jp} + 1 = 0$}
  {$e^{i\pi} + 1 = 0$}
\gnlaw{the rate of turning (Remark~\ref{rem:gn-euler-derivative})}
  {$\dfrac{d}{d\theta}\,e^{\theta\Jp} = \Jp\,e^{\theta\Jp}$}
  {$\dfrac{d}{d\theta}\,e^{i\theta} = i\,e^{i\theta}$}
\gnlaw{de Moivre's formula (\Cref{thm:gn-de-moivre})}
  {$(C_\theta + S_\theta\,\up)^n = C_{n\theta} + S_{n\theta}\,\up$}
  {$(\cos\theta + i\sin\theta)^n = {}$ \newline $\cos n\theta + i\sin n\theta$}
\gnlaw{the powers of the Composer Unit (\Cref{cor:gn-powers})}
  {$\Gp^n = |\Gp|^n C_{n\pb} + |\Gp|^{n-1}S_{n\pb}\,\Gp$ \newline $\quad = -|\Gp|^n\,\dfrac{\sin\,(n - 1)\pb}{\sin\pb} + |\Gp|^{n-1}\,\dfrac{\sin n\pb}{\sin\pb}\,\Gp$}
  {$i^n = \cos\dfrac{n\pi}{2} + i\sin\dfrac{n\pi}{2}$}
\gnlaw{the step rule}
  {$\Gp^{n+1} = 2|\Gp|\cos\pb\;\Gp^n - |\Gp|^2\,\Gp^{n-1}$}
  {$i^{n+1} = -i^{n-1}$}
\gnlaw{the numbers $\sin n\pb/\sin\pb$, for $n = 1, 2, 3, 4$}
  {$1$, \quad $2\cos\pb$, \quad $4\cos^2\pb - 1$, \quad $8\cos^3\pb - 4\cos\pb$}
  {$1$, \quad $0$, \quad $-1$, \quad $0$}
\gnlaw{the exponential of a General Number}
  {$e^z = e^{\Sh(z)}\,\bigl(C_{\Ht(z)} + S_{\Ht(z)}\,\up\bigr)$}
  {$e^{a + bi} = e^a(\cos b + i\sin b)$}
\gnlaw{the logarithm, for $-\pi < \theta \le \pi$}
  {$\ln z = \ln|z| + \theta\Jp = (\ln|z| - \theta\cot\pb) + \dfrac{\theta}{|\Gp|\sin\pb}\,\Gp$}
  {$\ln z = \ln r + i\theta$}
\end{gnlaws}
The numbers $\sin n\pb/\sin\pb$ are polynomials in $\cos\pb$, the Chebyshev polynomials $U_{n-1}(\cos\pb)$ of the second kind; the step rule makes each one from the two before it. The exponential of a General Number follows from $z = \Sh(z) + \Ht(z)\Jp$ and the rule $e^{w_1 + w_2} = e^{w_1}e^{w_2}$, and the logarithm from $z = |z|\,e^{\theta\Jp}$.

% keep the heading with its table: break the page here if fewer than 12 lines are left
\gnkeep{12}
%------------------------------------------------
\section{Roots}\label{app:f-roots}
%------------------------------------------------

\begin{gnlaws}
\gnlaw{the $n$-th roots of $z$ (\Cref{thm:gn-roots})}
  {$w_k = |z|^{1/n}\,\bigl(C_{\theta_k} + S_{\theta_k}\,\up\bigr)$, \quad $\theta_k = \dfrac{\theta + 2\pi k}{n}$, \quad $k = 0, \dots, n - 1$}
  {$w_k = |z|^{1/n}\,e^{i(\theta + 2\pi k)/n}$}
\gnlaw{the roots of unity}
  {$C_{2\pi k/n} + S_{2\pi k/n}\,\up$}
  {$e^{2\pi ik/n}$}
\gnlaw{the square roots of the Composer Unit (\Cref{cor:gn-sqrt})}
  {$\pm\dfrac{\sqrt{|\Gp|}}{2\cos(\pb/2)}\,(1 + \up)$}
  {$\pm\dfrac{1 + i}{\sqrt{2}}$}
\gnlaw{one step along each axis, squared}
  {$(1 + \up)^2 = 4\cos^2(\pb/2)\;\up$}
  {$(1 + i)^2 = 2i$}
\gnlaw{the quadratic equation $v^2 + Pv + Q = 0$}
  {$v = \dfrac{-P \pm \sqrt{P^2 - 4Q}}{2}$, with the square roots of the first line}
  {the same}
\end{gnlaws}
The fourth line follows from $\up^2 = 2\cos\pb\;\up - 1$ and $1 + \cos\pb = 2\cos^2(\pb/2)$. The fifth holds in every field in which each number has its two square roots, and the first line gives them.

% keep the heading with its table: break the page here if fewer than 22 lines are left
\gnkeep{22}
%------------------------------------------------
\section{The matrix of a General Number}\label{app:f-matrix}
%------------------------------------------------

\begin{gnlaws}
\gnlaw{the matrix of $z$}
  {$\begin{pmatrix} a & -|\Gp|^2 b \\ b & a + 2b|\Gp|\cos\pb \end{pmatrix}$}
  {$\begin{pmatrix} a & -b \\ b & a \end{pmatrix}$}
\gnlaw{products}
  {the matrix of $zw$ is the product of the matrices of $z$ and $w$}
  {the same}
\gnlaw{determinant and trace}
  {$\det = |z|^2$, \quad trace $= 2\Sh(z)$}
  {$\det = a^2 + b^2$, \quad trace $= 2a$}
\gnlaw{the matrix of the Composer Unit}
  {$\begin{pmatrix} 0 & -|\Gp|^2 \\ 1 & 2|\Gp|\cos\pb \end{pmatrix}$, which obeys the Rule of the Square: $M^2 = 2|\Gp|\cos\pb\,M - |\Gp|^2$}
  {$\begin{pmatrix} 0 & -1 \\ 1 & 0 \end{pmatrix}$, with $M^2 = -1$}
\end{gnlaws}
Multiplying by $z$ sends $1$ to $a + b\Gp$ and $\Gp$ to $-|\Gp|^2 b + (a + 2b|\Gp|\cos\pb)\,\Gp$, by Proposition~\ref{prop:un-real-product}; these two pairs of parts are the columns of the matrix.

% keep the heading with its table: break the page here if fewer than 16 lines are left
\gnkeep{16}
%------------------------------------------------
\section{The parts beside the shadow and the height}\label{app:f-parts}
%------------------------------------------------

A number has two pairs of coordinates: its parts, $\Re(z) = a$ and $\Com(z) = b\Gp$ (Definition~\ref{def:gn-parts}), counted in the leaning frame, and its shadow and height, counted in the upright frame (Proposition~\ref{prop:gn-shadow-height}). In each row the law for the parts comes first, and under it the same law for the shadow and the height; beside them stands the complex case, where the two pairs are one. Each row cites the statements that prove it, most of them in \Cref{subsec:gn-parts-beside}.

\begin{gnlawswide}
\gnlaw{the two pairs (Definitions~\ref{def:gn-parts} and~\ref{def:shadow}, Proposition~\ref{prop:gn-shadow-height})}
  {$\Re(z) = a$, \quad $\Com(z) = b\Gp$; \newline $\Sh(z) = a + b|\Gp|\cos\pb$, \quad $\Ht(z) = b|\Gp|\sin\pb$}
  {$\operatorname{Re} z = a$, \quad $\Com(z) = bi$; \newline the two pairs are one}
\gnlaw{from one pair to the other (Proposition~\ref{prop:gn-parts-readings})}
  {$\Re(z) = \Sh(z) - \Ht(z)\cot\pb$, \quad $\Com(z) = \Ht(z)\,\Gp/(|\Gp|\sin\pb)$; \newline $\Sh(z) = \Re(z) + \Sh\bigl(\Com(z)\bigr)$, \quad $\Ht(z) = \Ht\bigl(\Com(z)\bigr)$}
  {the shadow is the real part, and $\Com(z) = i\operatorname{Im} z$}
\gnlaw{sum and real multiple (Proposition~\ref{prop:gn-parts-readings})}
  {$\Re(z + w) = \Re(z) + \Re(w)$ and $\Re(kz) = k\,\Re(z)$ for every real $k$, and the same for $\Com$, for $\Sh$ and for $\Ht$}
  {the same for $\operatorname{Re}$ \newline and for $\operatorname{Im}$}
\gnlaw{magnitude (\Cref{thm:triangle}, Proposition~\ref{prop:gn-shadow-height})}
  {$|z|^2 = \Re(z)^2 + 2\cos\pb\;\Re(z)\,b|\Gp| + b^2|\Gp|^2$; \newline $|z|^2 = \Sh(z)^2 + \Ht(z)^2$}
  {$|z|^2 = a^2 + b^2$}
\gnlaw{from a number and its conjugate (Proposition~\ref{prop:gn-parts-readings})}
  {$\Re(z) = (\Gp\bar{z} - \overline{\Gp}z)/(\Gp - \overline{\Gp})$, \quad $\Com(z) = \Gp\,(z - \bar{z})/(\Gp - \overline{\Gp})$, \newline with $\Gp - \overline{\Gp} = 2|\Gp|\sin\pb\;\Jp$; \newline $\Sh(z) = (z + \bar{z})/2$, \quad $\Ht(z)\,\Jp = (z - \bar{z})/2$}
  {$\operatorname{Re} z = (z + \bar{z})/2$, \newline $i\operatorname{Im} z = (z - \bar{z})/2$}
\gnlaw{real numbers, pure Composer numbers and multiples of $\Jp$ (Proposition~\ref{prop:gn-parts-readings})}
  {$\Com(z) = 0$ exactly when $\bar{z} = z$, \newline $\Re(z) = 0$ exactly when $\Gp\bar{z} = \overline{\Gp}z$; \newline $\Ht(z) = 0$ exactly when $\bar{z} = z$, \newline $\Sh(z) = 0$ exactly when $\bar{z} = -z$}
  {real numbers: $\bar{z} = z$, \newline imaginary numbers: $\bar{z} = -z$}
\gnlaw{product (\Cref{thm:gn-parts-product})}
  {$\Re(zw) = \Re(z)\Re(w) - |\Gp|^2 bd$, \newline $\Com(zw) = \Sh(z)\Com(w) + \Sh(w)\Com(z)$; \newline $\Sh(zw) = \Sh(z)\Sh(w) - \Ht(z)\Ht(w)$, \newline $\Ht(zw) = \Sh(z)\Ht(w) + \Sh(w)\Ht(z)$}
  {$\operatorname{Re}(zw) = ac - bd$, \newline $\operatorname{Im}(zw) = ad + bc$}
\gnlaw{square (\Cref{thm:gn-parts-product})}
  {$\Re(z^2) = \Re(z)^2 - |\Gp|^2 b^2$, \quad $\Com(z^2) = 2\Sh(z)\Com(z)$; \newline $\Sh(z^2) = \Sh(z)^2 - \Ht(z)^2$, \quad $\Ht(z^2) = 2\Sh(z)\Ht(z)$}
  {$\operatorname{Re}(z^2) = a^2 - b^2$, \newline $\operatorname{Im}(z^2) = 2ab$}
\gnlaw{conjugate (Proposition~\ref{prop:gn-parts-conjugate})}
  {$\Re(\bar{z}) = 2\Sh(z) - \Re(z)$, \quad $\Com(\bar{z}) = -\Com(z)$, \newline so $\Sh(z) = \bigl[\Re(z) + \Re(\bar{z})\bigr]/2$; \newline $\Sh(\bar{z}) = \Sh(z)$, \quad $\Ht(\bar{z}) = -\Ht(z)$}
  {$\operatorname{Re}\bar{z} = a$, \newline $\operatorname{Im}\bar{z} = -b$}
\gnlaw{inverse (Proposition~\ref{prop:gn-parts-conjugate})}
  {$\Re(1/z) = \bigl[2\Sh(z) - \Re(z)\bigr]/|z|^2$, \quad $\Com(1/z) = -\Com(z)/|z|^2$; \newline $\Sh(1/z) = \Sh(z)/|z|^2$, \quad $\Ht(1/z) = -\Ht(z)/|z|^2$}
  {the two parts of \newline $\dfrac{a - bi}{a^2 + b^2}$}
\gnlaw{quotient, with the top $z\bar{w}$ (Proposition~\ref{prop:gn-parts-conjugate})}
  {$\Re(z/w) = \Re(z\bar{w})/|w|^2$ \ with \ $\Re(z\bar{w}) = \Re(z)\Re(\bar{w}) + |\Gp|^2 bd$, \newline $\Com(z/w) = \Com(z\bar{w})/|w|^2$ \ with \ $\Com(z\bar{w}) = \Sh(w)\Com(z) - \Sh(z)\Com(w)$; \newline $\Sh(z/w) = \Sh(z\bar{w})/|w|^2$ \ with \ $\Sh(z\bar{w}) = \Sh(z)\Sh(w) + \Ht(z)\Ht(w)$, \newline $\Ht(z/w) = \Ht(z\bar{w})/|w|^2$ \ with \ $\Ht(z\bar{w}) = \Sh(w)\Ht(z) - \Sh(z)\Ht(w)$}
  {the two parts of \newline $\dfrac{(ac + bd) + (bc - ad)\,i}{c^2 + d^2}$}
\gnlaw{a quarter turn, the product by $\Jp$ (\Cref{cor:gn-parts-turns})}
  {$\Re(\Jp z) = -\Ht(z) - \Sh(z)\cot\pb$, \quad $\Com(\Jp z) = \Sh(z)\,\Gp/(|\Gp|\sin\pb)$; \newline $\Sh(\Jp z) = -\Ht(z)$, \quad $\Ht(\Jp z) = \Sh(z)$}
  {$\operatorname{Re}(iz) = -b$, \newline $\operatorname{Im}(iz) = a$}
\gnlaw{the product by the Composer Unit (\Cref{cor:gn-parts-turns})}
  {$\Re(\Gp z) = -|\Gp|^2 b$, \quad $\Com(\Gp z) = \Re(\bar{z})\,\Gp$; \newline $\Sh(\Gp z) = |\Gp|\bigl(\Sh(z)\cos\pb - \Ht(z)\sin\pb\bigr)$, \newline $\Ht(\Gp z) = |\Gp|\bigl(\Sh(z)\sin\pb + \Ht(z)\cos\pb\bigr)$}
  {the product by $i$, \newline as in the line above}
\gnlaw{the powers of the Composer Unit (\Cref{thm:shadow-product}, Corollaries~\ref{cor:gn-powers} and~\ref{cor:gn-two-trig})}
  {$\Re(\Gp^n) = -|\Gp|^n\sin\,(n - 1)\pb/\sin\pb$, \quad $\Com(\Gp^n) = |\Gp|^{n-1}\,(\sin n\pb/\sin\pb)\,\Gp$; \newline $\Sh(\Gp^n) = |\Gp|^n\cos n\pb$, \quad $\Ht(\Gp^n) = |\Gp|^n\sin n\pb$}
  {$\operatorname{Re}(i^n) = \cos(n\pi/2)$, \newline $\operatorname{Im}(i^n) = \sin(n\pi/2)$}
\gnlaw{powers, and for $n = 1$ the polar form (\Cref{cor:gn-two-trig})}
  {$\Re(z^n) = |z|^n C_{n\theta}$, \quad $\Com(z^n) = |z|^n S_{n\theta}\,\up$; \newline $\Sh(z^n) = |z|^n\cos n\theta$, \quad $\Ht(z^n) = |z|^n\sin n\theta$}
  {$\operatorname{Re}(z^n) = r^n\cos n\theta$, \newline $\operatorname{Im}(z^n) = r^n\sin n\theta$}
\gnlaw{exponential (\Cref{cor:gn-two-trig})}
  {$\Re(e^z) = e^{\Sh(z)}\,C_{\Ht(z)}$, \quad $\Com(e^z) = e^{\Sh(z)}\,S_{\Ht(z)}\,\up$; \newline $\Sh(e^z) = e^{\Sh(z)}\cos\Ht(z)$, \quad $\Ht(e^z) = e^{\Sh(z)}\sin\Ht(z)$}
  {$\operatorname{Re}(e^z) = e^a\cos b$, \newline $\operatorname{Im}(e^z) = e^a\sin b$}
\gnlaw{signed area and dot product (Proposition~\ref{prop:gn-signed-area})}
  {$\Com(\bar{z}w) = (ad - bc)\,\Gp$, \quad $\Re(\bar{z}w) = \Re(\bar{z})\Re(w) + |\Gp|^2 bd$; \newline $\Ht(\bar{z}w) = \Sh(z)\Ht(w) - \Ht(z)\Sh(w)$, which is $(ad - bc)\,|\Gp|\sin\pb$, \newline $\Sh(\bar{z}w) = \Sh(z)\Sh(w) + \Ht(z)\Ht(w)$}
  {$\operatorname{Im}(\bar{z}w) = ad - bc$, \newline $\operatorname{Re}(\bar{z}w) = ac + bd$}
\gnlaw{size (Remark~\ref{rem:gn-long-parts})}
  {$|\Re(z)| \le |z|/\sin\pb$, \quad $|\Com(z)| \le |z|/\sin\pb$; \newline $|\Sh(z)| \le |z|$, \quad $|\Ht(z)| \le |z|$}
  {$|\operatorname{Re} z| \le |z|$, \newline $|\operatorname{Im} z| \le |z|$}
\end{gnlawswide}
From the third row on, the laws for the shadow and the height are the laws of the complex numbers, read in the upright frame, word for word; the first two rows pass from one frame to the other. The laws for the parts are the same arithmetic counted in the leaning frame, where the lean shows. Written out in the coordinates $a$, $b$, $c$ and $d$, two of them carry neither $\cos\pb$ nor $\sin\pb$, and nor do their special cases: the real part of $zw$, which is $ac - |\Gp|^2bd$, and the Composer part of $\bar{z}w$, which is $(ad - bc)\,\Gp$. The angle enters the first only through the length $|\Gp|$. In the row of real numbers, the lines for the shadow and the height follow from $z + \bar{z} = 2\Sh(z)$ and $z - \bar{z} = 2\Ht(z)\,\Jp$.

% keep the table of the six numbers together: break the page here if fewer than 14 lines are left
\gnkeep{14}
The parts, the shadow and the height of six numbers that recur:
\begin{center}
\small
\begin{tabular}{@{}lllll@{}}
\toprule
 & \spacedlowsmallcaps{real part} & \spacedlowsmallcaps{Composer part} & \spacedlowsmallcaps{shadow} & \spacedlowsmallcaps{height} \\
\midrule
$1$ & $1$ & $0$ & $1$ & $0$ \\
$\Gp$ & $0$ & $\Gp$ & $|\Gp|\cos\pb$ & $|\Gp|\sin\pb$ \\
$\overline{\Gp}$ & $2|\Gp|\cos\pb$ & $-\Gp$ & $|\Gp|\cos\pb$ & $-|\Gp|\sin\pb$ \\
$\Gp^2$ & $-|\Gp|^2$ & $2|\Gp|\cos\pb\;\Gp$ & $|\Gp|^2\cos 2\pb$ & $|\Gp|^2\sin 2\pb$ \\
$1/\Gp$ & $2\cos\pb/|\Gp|$ & $-\Gp/|\Gp|^2$ & $\cos\pb/|\Gp|$ & $-\sin\pb/|\Gp|$ \\
$\Jp$ & $-\cot\pb$ & $\Gp/(|\Gp|\sin\pb)$ & $0$ & $1$ \\
\bottomrule
\end{tabular}
\end{center}
At $\pb = \pi/2$ the six numbers are $1$, $i$, $-i$, $-1$, $-i$ and $i$. The first two columns then give the real part and $i$ times the imaginary part, and the last two the real part and the imaginary part.
